\section{Baseline model and analytical foundations}
\label{sec:toy}

The toy model isolates the interaction between population growth, market power in
current AI services, and the automation of AI research. Time is continuous. There is
one final good, a representative dynasty, competitive final-good firms, and one
integrated frontier developer. The developer monopolistically supplies AI services
and invests in the next generation of AI. The initial economy already uses AI
services in production; the model does not describe their first commercial adoption.

The model contains four linked blocks: labor allocation, current production, frontier
research, and goods allocation. The diagrams below present each block where its
equations are introduced. Variables repeated across diagrams connect the blocks.
Solid arrows denote within-period production or resource use; dashed arrows denote
state accumulation. This separation is important because the model is a dynamic
system with feedback, not a one-way production tree.

\tikzset{
    model-input/.style={draw, rounded corners=2pt, minimum width=2.7cm,
        minimum height=1.0cm, align=center, fill=gray!8},
    model-state/.style={draw, rounded corners=2pt, minimum width=2.7cm,
        minimum height=1.0cm, align=center, fill=blue!8},
    model-technology/.style={draw, rounded corners=2pt, minimum width=3.5cm,
        minimum height=1.0cm, align=center, fill=orange!9},
    model-allocation/.style={draw, rounded corners=2pt, minimum width=3.5cm,
        minimum height=1.0cm, align=center, fill=gray!8},
    model-flow/.style={->, semithick},
    model-dynamics/.style={->, semithick, dashed}
}

\subsection{Household}

Population grows at the exogenous rate $n$:
\begin{equation}
    \dot N_t=nN_t,
    \qquad N_0>0,\quad n\geq0.
    \label{eq:population}
\end{equation}
The representative dynasty has preferences
\begin{equation}
    \int_0^\infty e^{-\rho t}N_t\log c_t\dd t,
    \qquad
    c_t\equiv\frac{C_t}{N_t},
    \qquad \rho>n,
    \label{eq:utility}
\end{equation}
where $C_t$ is aggregate consumption. Each household member supplies one unit of
labor inelastically. Labor is allocated between final production, $L_t$, and
human AI research, $H_t$:
\begin{equation}
    L_t+H_t=N_t.
    \label{eq:labor-clearing}
\end{equation}
Figure~\ref{fig:labor-allocation-flow} shows the associated opportunity cost. An
additional worker assigned to AI research reduces the labor available for current
production.

\begin{figure}[H]
    \centering
    \begin{tikzpicture}[>=Latex,font=\small]
        \node[model-state]      (N)     at (-3.8, 0.0) {$N_t$\\population};
        \node[model-allocation] (labor) at ( 0.0, 0.0)
            {$N_t=L_t+H_t$\\labor allocation};
        \node[model-input]      (L)     at ( 3.8, 1.1) {$L_t$\\production labor};
        \node[model-input]      (H)     at ( 3.8,-1.1) {$H_t$\\human research};

        \draw[model-flow] (N) -- (labor);
        \draw[model-flow] (labor) -- (L);
        \draw[model-flow] (labor) -- (H);
    \end{tikzpicture}
    \caption{Allocation of population between current production and human AI
    research.}
    \label{fig:labor-allocation-flow}
\end{figure}

The dynasty owns the capital stock and the frontier developer. Capital evolves as
\begin{equation}
    \dot K_t=I_t-\delta K_t.
    \label{eq:capital-law}
\end{equation}
Let $k_t\equiv K_t/N_t$. The per-capita budget constraint is
\begin{equation}
    \dot k_t
    =
    (r_t-\delta-n)k_t+w_t
    +\frac{\pi^X_t-w_tH_t-\xi M_t}{N_t}
    -c_t,
    \label{eq:household-budget}
\end{equation}
where $\pi^X_t$ is the developer's operating profit before research expenditure.
The household's Euler equation is
\begin{equation}
    \frac{\dot c_t}{c_t}=r_t-\delta-\rho.
    \label{eq:euler}
\end{equation}

\subsection{Final-good firms}

Competitive final-good firms combine physical capital with a composite of production
labor and AI services:
\begin{align}
    Y_t
    &=
    K_t^\alpha Z_t^{1-\alpha},
    \qquad \alpha\in(0,1),
    \label{eq:final-production}\\
    Z_t
    &=
    \left[
        \omega_L L_t^{(\sigma_{XL}-1)/\sigma_{XL}}
        +\omega_X X_t^{(\sigma_{XL}-1)/\sigma_{XL}}
    \right]^{\sigma_{XL}/(\sigma_{XL}-1)}.
    \label{eq:service-composite}
\end{align}
The elasticity of substitution between capital and the service composite is one.
The elasticity between production labor and AI services is $\sigma_{XL}$. The
Cobb--Douglas composite is obtained by continuity when $\sigma_{XL}=1$. We assume
$\sigma_{XL}>0$. The pair $(\omega_L,\omega_X)$ contains positive CES distribution
weights satisfying $\omega_L+\omega_X=1$, not generally observed factor-income
shares; the latter are endogenous and vary with relative quantities when
$\sigma_{XL}\neq1$.

Define the contribution of AI services to the CES composite as
\begin{equation}
    s^X_t
    \equiv
    \frac{\omega_X X_t^{(\sigma_{XL}-1)/\sigma_{XL}}}
    {\omega_L L_t^{(\sigma_{XL}-1)/\sigma_{XL}}
        +\omega_X X_t^{(\sigma_{XL}-1)/\sigma_{XL}}}.
    \label{eq:ai-composite-share}
\end{equation}
Let $p^X_t$ denote the price of AI services. Profit maximization gives
\begin{align}
    r_t
    &=\alpha\frac{Y_t}{K_t},
    \label{eq:capital-demand}\\
    w_t
    &=(1-\alpha)(1-s^X_t)\frac{Y_t}{L_t},
    \label{eq:labor-demand}\\
    p^X_t
    &=(1-\alpha)s^X_t\frac{Y_t}{X_t},
    \label{eq:ai-demand}
\end{align}
and relative demand satisfies
\begin{equation}
    \frac{X_t}{L_t}
    =
    \left[
        \frac{\omega_X}{\omega_L}
        \frac{w_t}{p^X_t}
    \right]^{\sigma_{XL}}.
    \label{eq:relative-ai-demand}
\end{equation}

\subsection{Integrated AI monopolist}

The variable $U_t$ is a flow of \emph{inference compute}: standardized
computational resources used during period $t$ to run an existing model. It can be
measured in FLOPs, GPU-hours, or equivalent units of a reference processor per unit
of time; it is not a stock of chips. Let $\xi>0$ denote the common unit resource cost
of standardized compute. It summarizes the capital, energy, and other resources
needed to supply one unit of compute, independently of whether that unit is used for
inference or research.

The variable $X_t$ is a flow of quality-adjusted AI services used by final-good
firms. Its empirical counterpart could be benchmark-equivalent tokens, predictions
adjusted for accuracy, or equivalent task hours. This distinction follows the
measurement approach in \citet{korinekmckelvey2026}, who combine inference compute
with model capital and adjust inference output for quality, and is consistent with
the intelligence tiers documented by \citet{demireretal2025}.

The stock $A_t$ measures frontier AI capability. Under the normalization used here,
it also measures quality-adjusted services produced per unit of inference compute:
\begin{equation}
    X_t=A_tU_t.
    \label{eq:ai-services}
\end{equation}
Equation~\eqref{eq:ai-services} is an efficiency-units specification, not a
technological identity. Conditional on $A_t$, it imposes constant returns to
inference compute and compresses task coverage, quality, speed, and inference
efficiency into one scalar.

Figure~\ref{fig:current-production-flow} combines the AI-service technology with the
two nests used by final-good firms. It therefore shows how capability and inference
compute affect current output, holding research and goods allocation outside the
diagram.

\begin{figure}[H]
    \centering
    \resizebox{0.88\textwidth}{!}{%
    \begin{tikzpicture}[>=Latex,font=\small]
        \node[model-state]      (A) at (-8.0, 0.7) {$A_t$\\capability};
        \node[model-input]      (U) at (-8.0,-0.7) {$U_t$\\inference compute};
        \node[model-technology] (X) at (-4.5, 0.0) {$X_t=A_tU_t$\\AI services};
        \node[model-input]      (L) at (-4.5, 2.5) {$L_t$\\production labor};
        \node[model-technology] (Z) at ( 0.0, 1.2)
            {$Z_t=\operatorname{CES}(L_t,X_t)$};
        \node[model-state]      (K) at ( 0.0, 3.7) {$K_t$\\physical capital};
        \node[model-technology] (Y) at ( 4.0, 2.5)
            {$Y_t=K_t^\alpha Z_t^{1-\alpha}$};

        \draw[model-flow] (A) -- (X);
        \draw[model-flow] (U) -- (X);
        \draw[model-flow] (X) -- (Z);
        \draw[model-flow] (L) -- (Z);
        \draw[model-flow] (Z) -- (Y);
        \draw[model-flow] (K) -- (Y);
    \end{tikzpicture}
    }
    \caption{Current production. Capability and inference compute produce AI
    services; AI services and production labor form the service composite; the
    service composite and physical capital produce final output.}
    \label{fig:current-production-flow}
\end{figure}

\begin{remark}[Capability and inference efficiency]
The same state variable $A_t$ governs broad frontier capability and the productivity
of inference compute. This restriction keeps the service-market and research
mechanisms connected and makes the benchmark tractable, but it need not hold
quantitatively. A natural extension is $X_t=B(A_t)U_t$, with a separate mapping from
capability to inference efficiency. We retain $B(A)=A$ in the baseline and keep this
distinction as a robustness extension.
\end{remark}

The integrated developer recognizes the inverse demand generated by the final-good
sector and solves
\begin{equation}
    \max_{X_t\geq0}
    \left\{
        p^X_t(X_t;K_t,L_t)X_t
        -\xi\frac{X_t}{A_t}
    \right\}.
    \label{eq:monopoly-service-problem}
\end{equation}
Its operating profit is
\begin{equation}
    \pi^X_t=p^X_tX_t-\xi\frac{X_t}{A_t}.
    \label{eq:operating-profit}
\end{equation}

\subsection{Frontier research}

New AI capability is produced with effective research services, $E_t$:
\begin{align}
    \dot A_t
    &=
    \chi A_t^\phi E_t^\eta,
    \qquad \chi>0,\quad \phi<1,\quad \eta\in(0,1),
    \label{eq:ideas}\\
    E_t
    &=
    \left[
        \omega_H H_t^{(\sigma_{HM}-1)/\sigma_{HM}}
        +\omega_M M_t^{(\sigma_{HM}-1)/\sigma_{HM}}
    \right]^{\sigma_{HM}/(\sigma_{HM}-1)}.
    \label{eq:effective-research}
\end{align}
Here, $M_t$ is a flow of standardized automated-research services. It includes
training, post-training, experiments, synthetic-data generation, and inference
performed by AI agents engaged in R\&D. The baseline normalizes one unit of these
services to one unit in the research aggregator and assigns it the common resource
cost $\xi$. A capability-dependent conversion from compute into automated-research
services is left for an extension.

\subsubsection{Economic interpretation of the AI variables}

The model distinguishes the resources used to operate AI systems, the services those
systems produce, and the knowledge created by AI research. These objects are easy to
conflate because all of them may be described informally as ``AI.''
Table~\ref{tab:ai-variable-interpretation} states their economic interpretation.

\begin{table}[H]
    \centering
    \caption{Economic interpretation of the AI block}
    \label{tab:ai-variable-interpretation}
    \small
    \begin{tabular}{p{0.08\textwidth}p{0.16\textwidth}p{0.66\textwidth}}
        \toprule
        Variable & Type & Economic meaning \\
        \midrule
        $A_t$
        & \raggedright Intangible stock
        & Frontier AI capability. It is privately controlled by the developer but
        non-rival in use: applying it to one unit of inference compute does not reduce
        its availability for another. It raises the quantity and quality of AI services
        obtained from a unit of inference compute and, through $A_t^\phi$, affects the
        productivity of subsequent research. It is not the stock of chips, the amount
        of compute, or a count of trained models. In the baseline it compresses trained
        model capital, algorithms, and other knowledge relevant for AI performance into
        one state variable. \\
        \addlinespace
        $U_t$
        & \raggedright Current input flow
        & Inference compute used during period $t$ to operate existing AI systems. Its
        empirical counterpart is a flow of FLOPs, GPU-hours, or equivalent processor
        services. It is rival and has current resource cost $\xi U_t$. It does not
        directly add to the capability stock. \\
        \addlinespace
        $X_t$
        & \raggedright Service-output flow
        & Quality-adjusted AI services delivered to final-good firms. Possible units
        include benchmark-equivalent tokens, accuracy-adjusted predictions, or
        equivalent task-hours. It is not raw compute, raw tokens, or the developer's
        revenue; expenditure on these services is $p^X_tX_t$. \\
        \addlinespace
        $H_t$
        & \raggedright Human-input flow
        & Human labor assigned to improving frontier AI. It may be measured as
        researcher-hours or efficiency units of research labor. It excludes workers
        who operate data centers or market AI products but do not contribute directly
        to frontier research. Its opportunity cost is explicit in $L_t+H_t=N_t$, and
        its wage cost is $w_tH_t$. \\
        \addlinespace
        $M_t$
        & \raggedright Automated-input flow
        & Automated research services used to improve AI. They include training and
        post-training compute, experiments, synthetic-data generation, and inference
        performed by AI research agents. $M_t$ is not a stock of machines or trained
        models. The baseline expresses it in standardized compute-resource units, so
        its current resource cost is $\xi M_t$. \\
        \addlinespace
        $E_t$
        & \raggedright Research-output flow
        & Effective research services produced by combining $H_t$ and $M_t$. It is a
        CES index of economically productive research effort, not researcher headcount,
        physical compute, or an asset that can be stored. \\
        \addlinespace
        $\dot A_t$
        & \raggedright Innovation-output flow
        & The flow of newly produced capability per unit of time. It accumulates into
        $A_t$ but is not research expenditure and is not the proportional growth rate
        of capability; that growth rate is $\dot A_t/A_t$. Because the baseline has no
        capability depreciation, $\dot A_t$ is also the net addition to the stock. \\
        \bottomrule
    \end{tabular}
\end{table}

The production-accounting framework of \citet{korinekmckelvey2026} helps relate
these model variables to observable AI activity. In their framework, chips and
electricity produce compute, compute is allocated between inference and training,
model capital and inference compute produce quality-adjusted inference output, and
training compute and algorithms form new model capital. The approximate
correspondence with the present model is
\begin{equation}
    \begin{aligned}
        \text{inference compute} &\longleftrightarrow U_t,\\
        \text{quality-adjusted inference output} &\longleftrightarrow X_t,\\
        \text{training and R\&D compute} &\longleftrightarrow
        \text{part of }M_t,\\
        \text{model capital and algorithmic knowledge} &\longleftrightarrow
        \text{components of }A_t.
    \end{aligned}
    \label{eq:korinek-mckelvey-mapping}
\end{equation}
These correspondences are not identities. In particular, $M_t$ is broader than raw
training compute because it includes AI systems performing research tasks, while
$A_t$ is broader than trained-model capital because it also absorbs algorithmic and
other knowledge. Conversely, \citet{korinekmckelvey2026} model the upstream
production of compute from chips and electricity explicitly. The baseline here
collapses that sector into the common unit resource cost $\xi$.

The three equations in the AI block therefore describe different economic margins.
Equation~\eqref{eq:ai-services} converts a current input, $U_t$, into current services,
$X_t$, using the existing capability stock. Equation~\eqref{eq:effective-research}
aggregates human and automated inputs into effective research, $E_t$.
Equation~\eqref{eq:ideas} converts effective research into new capability, $\dot A_t$.
The first margin determines how intensively existing AI is used; the other two
determine how quickly the technology available in subsequent periods improves.

Figure~\ref{fig:research-flow} separates the production of effective research from
the law of motion for capability. Human researchers and automated research services
produce $E_t$. Effective research and existing capability determine $\dot A_t$, which
in turn accumulates into the capability stock.

\begin{figure}[H]
    \centering
    \begin{tikzpicture}[>=Latex,font=\small]
        \node[model-input]      (H)    at (-6.0, 1.0) {$H_t$\\human research};
        \node[model-input]      (M)    at (-6.0,-1.0) {$M_t$\\automated research};
        \node[model-technology] (E)    at (-2.5, 0.0)
            {$E_t=\operatorname{CES}(H_t,M_t)$};
        \node[model-state]      (A)    at (-2.5, 2.5) {$A_t$\\capability};
        \node[model-technology] (Adot) at ( 1.5, 1.25)
            {$\dot A_t=\chi A_t^\phi E_t^\eta$};

        \draw[model-flow] (H) -- (E);
        \draw[model-flow] (M) -- (E);
        \draw[model-flow] (E) -- (Adot);
        \draw[model-flow] (A) -- (Adot);
        \draw[model-dynamics] (Adot.east) -- (4.0,1.25) -- (4.0,3.7) --
            node[above] {capability accumulation} (-2.5,3.7) -- (A.north);
    \end{tikzpicture}
    \caption{Frontier research and capability accumulation.}
    \label{fig:research-flow}
\end{figure}

The elasticity $\sigma_{HM}>1$ governs substitution between human and automated
research. The pair $(\omega_H,\omega_M)$ contains positive CES distribution weights
satisfying $\omega_H+\omega_M=1$. The contribution of each input to effective research
is endogenous except in the Cobb--Douglas limit.
The restriction $\sigma_{HM}>1$ makes the two inputs gross substitutes and allows
positive research output when $H_t=0$. It does not impose that human research
disappears. That outcome depends on the equilibrium path of the wage relative to
$\xi$. The exponent $\phi$ captures how the existing stock of knowledge affects the
productivity of both research inputs.

\begin{assumption}[Stable baseline human-essential research benchmark]
\label{ass:stable-human-essential-benchmark}
Throughout the baseline analysis and the numerical experiments, except where
conditional result~\ref{cond:singularity-conditions} explicitly studies the boundary
and violation of the restriction, the parameters satisfy
\begin{equation}
    1-\phi
    -\eta\omega_M\frac{\omega_X}{\omega_L}
    >0.
    \label{eq:fully-cd-stability-condition}
\end{equation}
\end{assumption}

Assumption~\ref{ass:stable-human-essential-benchmark} ensures that diminishing
returns to capability dominate the output--research feedback when both production
and research are Cobb--Douglas. It does not impose the stricter AI-dominated
condition $1-\phi-\eta\omega_X/\omega_L>0$. It also does not rule out the separate
production-side acceleration mechanism generated by $\sigma_{XL}>1$.

The two CES elasticities operate in different nests. The production elasticity
$\sigma_{XL}$ governs substitution between $L_t$ and $X_t$. The research elasticity
$\sigma_{HM}$ governs substitution between $H_t$ and $M_t$. The variable $E_t$ is the
output of the research CES in \eqref{eq:effective-research}; it is not a separate
input combined with $H_t$.

\begin{remark}[Capability-dependent research productivity]
An extension replaces $M_t$ inside \eqref{eq:effective-research} with
$\lambda(A_t)M_t$. Its effective unit cost is then $\xi/\lambda(A_t)$, and the
automation condition depends on $w_t\lambda(A_t)/\xi$ rather than $w_t/\xi$. In the
AI-dominated limit, a power function $\lambda(A)=\lambda_0A^\psi$ changes the
capability exponent from $\phi$ to $\phi+\psi\eta$. The baseline sets
$\lambda(A)=1$ to isolate substitution between human and automated research from
capability-dependent improvements in research efficiency.
\end{remark}

Let $q_t$ be the developer's private shadow value of capability and let
$F(A_t,H_t,M_t)$ denote the right-hand side of \eqref{eq:ideas}. Its current-value
Hamiltonian is
\begin{equation}
    \mathcal H^P_t
    =
    \pi^X_t-w_tH_t-\xi M_t
    +q_tF(A_t,H_t,M_t).
    \label{eq:private-hamiltonian}
\end{equation}
More explicitly, taking the paths of $r_t$, $w_t$, $K_t$, and $L_t$ as given, the
developer's global problem is
\begin{equation}
\begin{aligned}
    &\max_{\{X_t,H_t,M_t,A_t\}_{t\geq0}}
    \int_0^\infty
    \exp\!\left[-\int_0^t(r_s-\delta)\dd s\right]
    \left[\pi^X_t-w_tH_t-\xi M_t\right]\dd t,\\
    &\text{subject to }\dot A_t=F(A_t,H_t,M_t),
    \qquad A_0\text{ given}.
\end{aligned}
    \label{eq:developer-market-value-program}
\end{equation}
The pointwise inverse demand faced by the developer is the one induced by final-good
firms. Writing the global program explicitly matters because the canonical first-order
and transversality conditions below need not by themselves be sufficient for a global
optimum when the state--control problem is nonconcave.

\subsection{Social planner}

Let $Q_t$ denote the social shadow value of capability. Unlike $q_t$, $Q_t$ includes
consumer surplus from a larger supply of AI services and knowledge benefits that the
developer cannot appropriate. The comparison between $Q_t$ and $q_t$ isolates the
knowledge and appropriability wedge.

\subsection{Equilibrium}

A decentralized equilibrium is a sequence
\[
    \{C_t,c_t,I_t,K_t,N_t,L_t,H_t,X_t,U_t,M_t,A_t\}_{t\geq0}
\]
and prices
\[
    \{r_t,w_t,p^X_t,q_t\}_{t\geq0}
\]
such that:
\begin{enumerate}[label=(\roman*)]
    \item the representative dynasty maximizes \eqref{eq:utility}, taking prices and
    developer profits as given;
    \item final-good firms maximize profits under
    \eqref{eq:final-production}--\eqref{eq:service-composite};
    \item the integrated developer chooses current AI services and frontier research
    to maximize its market value;
    \item population, capital, and capability follow
    \eqref{eq:population}, \eqref{eq:capital-law}, and \eqref{eq:ideas};
    \item labor satisfies \eqref{eq:labor-clearing}; and
    \item the aggregate resource constraint is
    \begin{equation}
        C_t+I_t+\xi(U_t+M_t)=Y_t.
        \label{eq:resource-constraint}
    \end{equation}
\end{enumerate}

Figure~\ref{fig:goods-allocation-flow} isolates the final resource constraint.
Inference and automated research use the same standardized resource cost, $\xi$.
Investment is the only current use of output that accumulates into physical capital.

\begin{figure}[H]
    \centering
    \begin{tikzpicture}[>=Latex,font=\small]
        \node[model-technology] (Y)     at (-5.8, 0.0) {$Y_t$\\final output};
        \node[model-allocation, minimum width=4.1cm] (goods) at (-1.5,0.0)
            {$Y_t=C_t+I_t+\xi(U_t+M_t)$\\goods allocation};
        \node[model-input]      (C)     at ( 2.5, 2.4) {$C_t$\\consumption};
        \node[model-input]      (I)     at ( 2.5, 0.8) {$I_t$\\investment};
        \node[model-input]      (U)     at ( 2.5,-0.8) {$U_t$\\inference compute};
        \node[model-input]      (M)     at ( 2.5,-2.4) {$M_t$\\automated research};
        \node[model-state]      (K)     at ( 7.0, 0.8) {$K_t$\\physical capital};

        \draw[model-flow] (Y) -- (goods);
        \draw[model-flow] (goods) -- (C);
        \draw[model-flow] (goods) -- (I);
        \draw[model-flow] (goods) -- node[below left] {$\xi U_t$} (U);
        \draw[model-flow] (goods) -- node[below] {$\xi M_t$} (M);
        \draw[model-dynamics] (I) -- (K);
    \end{tikzpicture}
    \caption{Allocation of final output and physical-capital accumulation.}
    \label{fig:goods-allocation-flow}
\end{figure}

The following system records the market-clearing identities and canonical conditions
used to compute an interior equilibrium candidate. On any interval over which
the allocation is interior, all quantities and the private capability shadow value are positive,
$0<H_t<N_t$, and $0<\varepsilon^X_t<1$. An equilibrium trajectory satisfies
\begin{subequations}
\label{eq:equilibrium-trajectory-system}
\begin{align}
    \dot N_t&=nN_t,
    \label{eq:eqsys-population}\\
    \dot K_t&=I_t-\delta K_t,
    \label{eq:eqsys-capital}\\
    \dot K_t&=(r_t-\delta)K_t+w_tN_t+\pi^X_t-w_tH_t-\xi M_t-C_t,
    \label{eq:eqsys-household-budget}\\
    \frac{\dot C_t}{C_t}&=n+r_t-\delta-\rho,
    \label{eq:eqsys-consumption}\\
    c_t&=\frac{C_t}{N_t},
    \label{eq:eqsys-per-capita-consumption}\\
    C_t+I_t+\xi(U_t+M_t)&=Y_t,
    \label{eq:eqsys-resources}\\
    L_t+H_t&=N_t,
    \label{eq:eqsys-labor}\\
    Y_t&=K_t^\alpha Z_t^{1-\alpha},
    \label{eq:eqsys-output}\\
    Z_t&=
    \left[
        \omega_L L_t^{(\sigma_{XL}-1)/\sigma_{XL}}
        +\omega_X X_t^{(\sigma_{XL}-1)/\sigma_{XL}}
    \right]^{\sigma_{XL}/(\sigma_{XL}-1)},
    \label{eq:eqsys-production}\\
    X_t&=A_tU_t,
    \label{eq:eqsys-ai-services}\\
    \dot A_t=F(A_t,H_t,M_t)&=\chi A_t^\phi E_t^\eta,
    \label{eq:eqsys-capability}\\
    E_t&=
    \left[
        \omega_H H_t^{(\sigma_{HM}-1)/\sigma_{HM}}
        +\omega_M M_t^{(\sigma_{HM}-1)/\sigma_{HM}}
    \right]^{\sigma_{HM}/(\sigma_{HM}-1)},
    \label{eq:eqsys-research}\\
    E_{H,t}&=\omega_H\left(\frac{E_t}{H_t}\right)^{1/\sigma_{HM}},
    \label{eq:eqsys-human-research-marginal-product}\\
    E_{M,t}&=\omega_M\left(\frac{E_t}{M_t}\right)^{1/\sigma_{HM}},
    \label{eq:eqsys-machine-research-marginal-product}\\
    s^X_t&=
    \frac{\omega_X X_t^{(\sigma_{XL}-1)/\sigma_{XL}}}
    {\omega_L L_t^{(\sigma_{XL}-1)/\sigma_{XL}}
        +\omega_X X_t^{(\sigma_{XL}-1)/\sigma_{XL}}},
    \label{eq:eqsys-ai-share}\\
    \varepsilon^X_t&=
    \frac{1-s^X_t}{\sigma_{XL}}+\alpha s^X_t,
    \label{eq:eqsys-inverse-elasticity}\\
    r_t&=\alpha\frac{Y_t}{K_t},
    \label{eq:eqsys-capital-price}\\
    w_t&=(1-\alpha)(1-s^X_t)\frac{Y_t}{L_t},
    \label{eq:eqsys-wage}\\
    p^X_t&=(1-\alpha)s^X_t\frac{Y_t}{X_t},
    \label{eq:eqsys-ai-price}\\
    Y_t&=r_tK_t+w_tL_t+p^X_tX_t,
    \label{eq:eqsys-final-firm-zero-profit}
\end{align}
\end{subequations}
\begin{subequations}
\label{eq:equilibrium-optimality-system}
\begin{align}
    X_t&\in
    \underset{\widetilde X\geq0}{\arg\max}\,
    \left\{
        p^X_t(\widetilde X;K_t,L_t)\widetilde X
        -\frac{\xi}{A_t}\widetilde X
    \right\},
    \label{eq:eqsys-monopoly-global}\\
    p^X_t(1-\varepsilon^X_t)&=\frac{\xi}{A_t},
    \label{eq:eqsys-monopoly-foc}\\
    -\frac{X_t}{p^X_t}
    \frac{\partial}{\partial X_t}
    \left[p^X_t(1-\varepsilon^X_t)\right]
    &=
    \varepsilon^X_t(1-\varepsilon^X_t)
    +\left(\alpha-\frac{1}{\sigma_{XL}}\right)
    \left(1-\frac{1}{\sigma_{XL}}\right)s^X_t(1-s^X_t)>0,
    \label{eq:eqsys-monopoly-soc}\\
    \pi^X_t&=p^X_tX_t-\xi U_t,
    \label{eq:eqsys-profit}\\
    w_t&=q_tF_{H,t},
    \label{eq:eqsys-human-research-foc}\\
    \xi&=q_tF_{M,t},
    \label{eq:eqsys-machine-research-foc}\\
    F_{H,t}&=\chi\eta A_t^\phi E_t^{\eta-1}E_{H,t},
    \label{eq:eqsys-human-research-product}\\
    F_{M,t}&=\chi\eta A_t^\phi E_t^{\eta-1}E_{M,t},
    \label{eq:eqsys-machine-research-product}\\
    \dot q_t&=(r_t-\delta)q_t-\pi^X_{A,t}-q_tF_{A,t},
    \label{eq:eqsys-costate}\\
    \pi^X_{A,t}
    \equiv
    \left.\frac{\partial\pi^X_t}{\partial A_t}\right|_{X_t,K_t,L_t}
    &=\frac{\xi X_t}{A_t^2},
    \label{eq:eqsys-profit-derivative}\\
    F_{A,t}&=\phi\frac{\dot A_t}{A_t}.
    \label{eq:eqsys-capability-derivatives}
\end{align}
\end{subequations}
The global condition~\eqref{eq:eqsys-monopoly-global} is part of the equilibrium
system. The first-order condition~\eqref{eq:eqsys-monopoly-foc} alone would not be
sufficient if the monopolist's objective had multiple stationary points. Equation
\eqref{eq:eqsys-monopoly-soc} verifies a strict local maximum and provides a direct
numerical diagnostic. Given the maintained concavity of research production,
$\eta\in(0,1)$, the research first-order conditions characterize the developer's
optimal input scale and mix whenever both inputs are positive.

Systems~\eqref{eq:equilibrium-trajectory-system}--
\eqref{eq:equilibrium-optimality-system}, the initial conditions, and the
transversality or singular-endpoint conditions are therefore necessary for an
interior equilibrium. They are sufficient for the household, the competitive
final-good firms, and the developer's pointwise control choices. A full decentralized
equilibrium additionally requires the developer's path to solve the global program
\eqref{eq:developer-market-value-program}. The usual concavity sufficiency theorem
applies, for example, under restrictions that make the Hamiltonian jointly concave in
the state and controls. The baseline restrictions do not impose such global
concavity, so satisfying the canonical system is not, by itself, an existence,
uniqueness, or global-optimality proof.

The initial stocks $N_0$, $K_0$, and $A_0$ are given. Aggregate consumption $C_0$ and the private
shadow value $q_0$ jump to place the economy on a perfect-foresight path satisfying
the household and developer transversality conditions:
\begin{align}
    \lim_{t\to\infty}
    e^{-(\rho-n)t}\frac{K_t}{C_t}&=0,
    \label{eq:household-tvc}\\
    \lim_{t\to\infty}
    \exp\!\left[-\int_0^t(r_s-\delta)\dd s\right]q_tA_t&=0.
    \label{eq:developer-tvc}
\end{align}
If a control reaches a boundary, its equality in
\eqref{eq:eqsys-monopoly-foc} or
\eqref{eq:eqsys-human-research-foc}--\eqref{eq:eqsys-machine-research-foc} is
replaced by the corresponding Kuhn--Tucker condition. For the conditional finite-singularity paths
studied below, the numerical solution instead uses the limiting endpoint conditions
$K_t/C_t\to0$ and $q_tA_t/C_t\to0$ as $t\uparrow T^*$, together with the remaining
asymptotic restrictions in conditional
result~\ref{cond:high-sigma-singular-scaling}.

There is no common physical-capacity constraint between $U_t$ and $M_t$. The two
uses of compute compete through the final-good resource constraint, not through an
additional exogenous cap.

\subsection{Analytical foundations for the regime analysis}
\label{sec:analytical-results}

This subsection derives the analytical results used to distinguish the growth
regimes. Propositions report results established from the model's equations
without assuming a path for an endogenous variable. Conditional results instead
characterize necessary limits or the implications of an explicitly stated
closure. The proofs of the propositions and the derivations of the conditional
results are collected in the appendix. Section~\ref{sec:numerical} then returns to
each regime separately and places its analytical characterization next to an
equilibrium simulation. The simulations show how the relevant mechanisms interact
during the transition; they are not used to prove the analytical claims.

The analytical map has two technological dimensions. Research is either
human-essential ($\sigma_{HM}=1$) or permits human and automated inputs to become
gross substitutes ($\sigma_{HM}>1$). Final production makes labor and AI services
gross complements, Cobb--Douglas inputs, or gross substitutes according as
$\sigma_{XL}<1$, $\sigma_{XL}=1$, or $\sigma_{XL}>1$.
Table~\ref{tab:analytical-regime-map} organizes the resulting six cases. These are
technology classes, not six unique equilibrium outcomes. The sign of the recursive
feedback denominator is a separate dimension: it determines whether diminishing
returns stabilize capability growth within a given technology class.

\begin{table}[H]
\centering
\small
\caption{Six technology classes and their analytical status}
\label{tab:analytical-regime-map}
\begin{tabular}{@{}P{0.13\textwidth}P{0.25\textwidth}P{0.25\textwidth}P{0.25\textwidth}@{}}
\toprule
Research nest
& $\sigma_{XL}<1$
& $\sigma_{XL}=1$
& $\sigma_{XL}>1$ \\
\midrule
$\sigma_{HM}=1$
& Any stated interior constant-growth limit has zero per-capita output growth.
& $D_{\mathrm{CD}}$ governs the reduced recursive-feedback dynamics.
& Positive finite capability growth is incompatible with an interior finite-growth steady state. \\
\addlinespace
$\sigma_{HM}>1$
& The same necessary zero-growth limit applies; the research mix remains endogenous.
& The research-feedback denominator changes with the automated contribution.
& The same nonexistence result applies; an AI-dominated singular branch is characterized conditionally. \\
\bottomrule
\end{tabular}

\vspace{0.4em}
\begin{minipage}{0.95\textwidth}
\footnotesize
\textit{Notes:} The low-elasticity entries report necessary rates for the stated
nondegenerate limit, not existence theorems. The high-elasticity entries condition on
positive constant or unbounded capability and interior allocations. When both nests
are Cobb--Douglas, $D_{\mathrm{CD}}$ is constant. With CES research, the feedback
denominator changes with the automated contribution. When $\sigma_{XL}\neq1$, the
production share also changes, so $D_{\mathrm{CD}}$ is a benchmark rather than a
global classifier.
\end{minipage}
\end{table}

The six cells do not require six independent aggregate-growth analyses. The
$\sigma_{XL}<1$ necessary-rate result and the $\sigma_{XL}>1$ nonexistence result
apply under either research technology. In those two production regimes,
$\sigma_{HM}$ changes research composition and transition speed rather than the
headline classification. At $\sigma_{XL}=1$, by contrast, human essentiality
determines whether recursive feedback is governed by the constant
$D_{\mathrm{CD}}$ or by the endogenous path $D(s^M_t)$. The analytical discussion
therefore has four substantive regimes: production complementarity, human-essential
Cobb--Douglas growth, substitutable-research Cobb--Douglas growth, and production
gross substitution.

\subsubsection{A capital-deepening benchmark}
\label{sec:capital-deepening-benchmark}

Before analyzing AI services, consider a counterfactual in which the substitute for
labor is an ordinary reproducible capital service. Let total capital be divided
between an outer production input and a second input inside the CES composite:
\begin{align}
    \widetilde Y_t
    &=
    \widetilde K_{1t}^{\alpha}\widetilde Z_t^{1-\alpha},
    \label{eq:capital-benchmark-production}\\
    \widetilde Z_t
    &=
    \left[
        \omega_L L_t^{\varrho_K}
        +\omega_K\widetilde K_{2t}^{\varrho_K}
    \right]^{1/\varrho_K},
    \qquad
    \varrho_K\equiv
    \frac{\sigma_{KL}-1}{\sigma_{KL}},
    \qquad \omega_L+\omega_K=1,
    \qquad \omega_L,\omega_K\in(0,1),
    \label{eq:capital-benchmark-ces}\\
    \widetilde K_{1t}
    &=\zeta\widetilde K_t,
    \qquad
    \widetilde K_{2t}=(1-\zeta)\widetilde K_t,
    \qquad \zeta\in(0,1),
    \label{eq:capital-benchmark-allocation}\\
    \dot{\widetilde K}_t
    &=s_K\widetilde Y_t-\delta\widetilde K_t,
    \qquad s_K\in(0,1),
    \qquad
    \widetilde C_t=(1-s_K)\widetilde Y_t,
    \qquad g_{L,t}=n.
    \label{eq:capital-benchmark-accumulation}
\end{align}
The fixed allocation in \eqref{eq:capital-benchmark-allocation} is an explicit
closure used to isolate capital deepening. It is not imposed on the baseline AI
economy. Equivalently,
$\widetilde K_{1t}/\widetilde K_{2t}=\zeta/(1-\zeta)$ is an arbitrary positive
constant.

\begin{conditionalresult}[Capital deepening does not generate accelerating growth;
\citealp{pitchford1960}; \citealp{duffypapageorgiou2000}]
\label{cond:capital-deepening-benchmark}
Suppose $\sigma_{KL}>1$ and define
\begin{equation}
    B_K(\zeta)
    \equiv
    \zeta^\alpha(1-\zeta)^{1-\alpha}
    \omega_K^{(1-\alpha)/\varrho_K}.
    \label{eq:capital-benchmark-ak-coefficient}
\end{equation}
Then $\widetilde Y_t/\widetilde K_t\to B_K(\zeta)$ as
$\widetilde K_t/L_t\to\infty$. If
$s_KB_K(\zeta)<\delta+n$, the capital--labor ratio converges to a unique finite
steady state. If $s_KB_K(\zeta)>\delta+n$, the economy converges instead to an
asymptotic balanced-growth path with
\begin{align}
    g_{\widetilde K}
    =g_{\widetilde Y}
    &=s_KB_K(\zeta)-\delta,
    \label{eq:capital-benchmark-aggregate-growth}\\
    g_{\widetilde Y/L}
    &=s_KB_K(\zeta)-\delta-n>0.
    \label{eq:capital-benchmark-per-capita-growth}
\end{align}
On this path, labor's share converges to zero while the wage grows at
$g_{\widetilde Y/L}/\sigma_{KL}$. At the equality
$s_KB_K(\zeta)=\delta+n$, the capital--labor ratio diverges but its growth rate
converges to zero. In none of the three cases does the output growth rate diverge.
\end{conditionalresult}

The fixed capital composition is generally not the allocation chosen by a firm that
can move capital freely between the two uses. Let
\begin{equation}
    s^{K_2}_t
    \equiv
    \frac{\omega_K\widetilde K_{2t}^{\varrho_K}}
    {\omega_LL_t^{\varrho_K}
    +\omega_K\widetilde K_{2t}^{\varrho_K}}
    \label{eq:capital-benchmark-inner-share}
\end{equation}
denote the contribution of the second capital service inside the CES composite.
Equalizing the marginal returns to the two uses requires
\begin{equation}
    \frac{\widetilde K_{1t}}{\widetilde K_{2t}}
    =
    \frac{\alpha}{(1-\alpha)s^{K_2}_t}.
    \label{eq:capital-benchmark-optimal-allocation}
\end{equation}
Because $s^{K_2}_t$ generally changes with the capital--labor ratio, the optimal
capital composition generally changes over time. Under gross substitution,
$s^{K_2}_t\to1$ and the ratio converges to $\alpha/(1-\alpha)$. Allowing this
reallocation changes the asymptotic coefficient but not the benchmark's central
result: ordinary accumulable capital produces an asymptotically linear technology,
not an increasing growth rate.

This comparison separates factor substitution from recursive innovation. Replacing
labor with reproducible capital can eliminate labor as an asymptotic constraint, but
it leaves $\widetilde Y/\widetilde K$ finite. In the baseline model, $X_t=A_tU_t$
and capability is itself produced with human and automated research. Capability can
therefore raise current output and the productivity of the activity that produces
further capability. Any acceleration beyond the $AK$ benchmark must come from that
additional feedback rather than from gross substitution alone. The benchmark also
shows that a declining labor share together with a rising wage is not specific to
AI; the distinctive object in this paper is the interaction of that distributional
outcome with recursive research and market power.

\subsubsection{Wages and the labor share under an AI expansion}

Let
\begin{equation}
    \ell^Y_t\equiv\frac{w_tL_t}{Y_t}
    =(1-\alpha)(1-s^X_t)
    \label{eq:production-labor-share}
\end{equation}
denote production labor's share of final output. The wage and this share need not
move in the same direction because the wage depends on both the production-labor
share and output per production worker.

\begin{proposition}[\citealp{autorkausik2026}; \citealp{yango2026}]
\label{prop:wage-share-thresholds}
Holding $K_t$ and $L_t$ fixed, an expansion of AI services satisfies
\begin{align}
    \frac{\partial\log Y_t}{\partial\log X_t}
    &=(1-\alpha)s^X_t>0,
    \label{eq:output-ai-elasticity}\\
    \frac{\partial\log \ell^Y_t}{\partial\log X_t}
    &=-\left(1-\frac{1}{\sigma_{XL}}\right)s^X_t,
    \label{eq:labor-share-ai-elasticity}\\
    \frac{\partial\log w_t}{\partial\log X_t}
    &=\left(\frac{1}{\sigma_{XL}}-\alpha\right)s^X_t.
    \label{eq:wage-ai-elasticity}
\end{align}
Consequently, the production-labor share falls if and only if $\sigma_{XL}>1$, whereas
the wage rises if and only if $\sigma_{XL}<1/\alpha$. Since $\alpha\in(0,1)$, there is
a nonempty interval
\begin{equation}
    1<\sigma_{XL}<\frac{1}{\alpha}
    \label{eq:wage-share-decoupling-region}
\end{equation}
in which an expansion of AI raises the wage while reducing production labor's share
of output.
\end{proposition}

Because researchers and production workers receive the same wage, define the
aggregate labor-income share and the human-research employment share as
\begin{equation}
    \ell^N_t
    \equiv
    \frac{w_tN_t}{Y_t}
    =
    \frac{N_t}{L_t}\ell^Y_t
    =
    \frac{\ell^Y_t}{1-h_t},
    \qquad
    h_t\equiv\frac{H_t}{N_t}.
    \label{eq:aggregate-labor-share}
\end{equation}
The production and aggregate labor-income shares coincide only when no workers are
employed in research.

\begin{proposition}[Dynamic labor-income-share conditions]
\label{prop:dynamic-labor-share}
Along an equilibrium transition, the two labor-income shares obey
\begin{align}
    g_{\ell^Y,t}
    &=
    -\left(1-\frac{1}{\sigma_{XL}}\right)s^X_t
    \left(g_{X,t}-g_{L,t}\right),
    \label{eq:dynamic-production-labor-share}\\
    g_{\ell^N,t}
    &=
    -\left(1-\frac{1}{\sigma_{XL}}\right)s^X_t
    \left(g_{X,t}-g_{L,t}\right)
    +
    \frac{\dot h_t}{1-h_t}
    \notag\\
    &=
    g_{\ell^Y,t}+n-g_{L,t}.
    \label{eq:dynamic-aggregate-labor-share}
\end{align}
The production-labor share rises if
\begin{equation}
    \left(1-\frac{1}{\sigma_{XL}}\right)
    \left(g_{X,t}-g_{L,t}\right)<0
    \label{eq:production-labor-share-rise-condition}
\end{equation}
and falls if the inequality is reversed. The aggregate labor-income share rises if
\begin{equation}
    \frac{\dot h_t}{1-h_t}
    >
    \left(1-\frac{1}{\sigma_{XL}}\right)s^X_t
    \left(g_{X,t}-g_{L,t}\right),
    \label{eq:aggregate-labor-share-rise-condition}
\end{equation}
and falls if the inequality is reversed.
\end{proposition}

If AI services grow faster than production labor, the production-labor share rises
when $\sigma_{XL}<1$, is constant when $\sigma_{XL}=1$, and falls when $\sigma_{XL}>1$. The
aggregate share has an additional reallocation term. An increasing human-research
share raises aggregate labor income relative to the production-labor share, while a
declining human-research share lowers it. If $h_t$ is constant, the two shares have
the same growth rate and the familiar threshold $\sigma_{XL}=1$ applies to both.

Equation~\eqref{eq:wage-ai-elasticity} decomposes the wage response into a
productivity effect and a displacement effect:
\begin{equation}
    \frac{\partial\log w_t}{\partial\log X_t}
    =
    \underbrace{(1-\alpha)s^X_t}_{\text{productivity effect}}
    -
    \underbrace{\left(1-\frac{1}{\sigma_{XL}}\right)s^X_t}
    _{\text{displacement effect}}.
    \label{eq:wage-decomposition}
\end{equation}
The first term is the increase in output per production worker. The second is the
decline in labor's income share. This decomposition maps the nested production
function into the productivity and displacement channels emphasized by
\citet{acemoglurestrepo2019}.

The qualitative decoupling between the wage and labor's share is the central result
of \citet[Theorem~1]{autorkausik2026}, who establish it for a general competitive
constant-returns technology. For a nested Cobb--Douglas--CES technology,
\citet[Lemma~4.1 and Proposition~4.2]{yango2026} derive the exact elasticities and
the two thresholds in Proposition~\ref{prop:wage-share-thresholds}. We use this
comparative static as a benchmark for the model's distinct dynamic and
market-structure mechanisms.

Allowing all quantities to change, the instantaneous wage growth rate obeys the
identity
\begin{equation}
    g_{w,t}
    =
    \alpha\left(g_{K,t}-g_{L,t}\right)
    +
    s^X_t\left(\frac{1}{\sigma_{XL}}-\alpha\right)
    \left(g_{X,t}-g_{L,t}\right).
    \label{eq:dynamic-wage-decomposition}
\end{equation}
When the research-employment share is constant, $g_{L,t}=n$. Capital deepening can
therefore raise wages even when the direct expansion of AI services lowers them.
Conversely, $\sigma_{XL}<1/\alpha$ is not sufficient for wage growth along an equilibrium
transition if capital per production worker contracts. The monopoly problem below
determines the supply of $X_t$, while household saving and the resource constraint
determine $K_t$; the equilibrium wage response depends on both paths.

\subsubsection{Monopoly supply of AI services}

Holding $K_t$ and $L_t$ fixed, the absolute elasticity of inverse demand faced by
the monopolist is
\begin{equation}
    \varepsilon^X_t
    \equiv
    -\frac{\partial\log p^X_t}{\partial\log X_t}
    =
    \frac{1-s^X_t}{\sigma_{XL}}+\alpha s^X_t.
    \label{eq:inverse-demand-elasticity}
\end{equation}
For an interior solution with $\varepsilon^X_t<1$, the first-order condition is
\begin{equation}
    p^X_t(1-\varepsilon^X_t)=\frac{\xi}{A_t}.
    \label{eq:monopoly-ai-price}
\end{equation}
The markup and operating rent are
\begin{align}
    \mu^X_t
    &\equiv
    \frac{p^X_t}{\xi/A_t}
    =
    \frac{1}{1-\varepsilon^X_t}>1,
    \label{eq:monopoly-markup}\\
    \pi^X_t
    &=
    \varepsilon^X_t p^X_tX_t.
    \label{eq:monopoly-rents}
\end{align}
Conditional on $A_t$, $K_t$, and $L_t$, the monopolist supplies less AI than an
allocation that prices services at marginal cost. The household receives the rents,
so the toy model retains the output distortion and the innovation incentive from
market power while abstracting from concentrated ownership.

\subsubsection{Research input choice and gradual automation}

The marginal products of the two research inputs in
\eqref{eq:effective-research} are
\begin{equation}
    E_{H,t}
    =\omega_H\left(\frac{E_t}{H_t}\right)^{1/\sigma_{HM}},
    \qquad
    E_{M,t}
    =\omega_M\left(\frac{E_t}{M_t}\right)^{1/\sigma_{HM}}.
    \label{eq:research-ces-marginal-products}
\end{equation}
The developer's first-order conditions are therefore
\begin{align}
    w_t
    &=
    q_t\chi\eta A_t^\phi E_t^{\eta-1}
    \omega_H\left(\frac{E_t}{H_t}\right)^{1/\sigma_{HM}},
    \label{eq:foc-human}\\
    \xi
    &=
    q_t\chi\eta A_t^\phi E_t^{\eta-1}
    \omega_M\left(\frac{E_t}{M_t}\right)^{1/\sigma_{HM}}.
    \label{eq:private-compute-foc}
\end{align}
Their ratio determines the research-input mix:
\begin{equation}
    \frac{M_t}{H_t}
    =
    \left[
        \frac{\omega_M}{\omega_H}\frac{w_t}{\xi}
    \right]^{\sigma_{HM}}.
    \label{eq:research-input-ratio}
\end{equation}

The unit cost of effective research services is
\begin{equation}
    P^E_t
    =
    \left[
        \omega_H^{\sigma_{HM}} w_t^{1-\sigma_{HM}}
        +\omega_M^{\sigma_{HM}}\xi^{1-\sigma_{HM}}
    \right]^{1/(1-\sigma_{HM})}.
    \label{eq:research-unit-cost}
\end{equation}
Conditional input demands are
\begin{equation}
    H_t=\omega_H^{\sigma_{HM}}
    \left(\frac{P^E_t}{w_t}\right)^{\sigma_{HM}} E_t,
    \qquad
    M_t=\omega_M^{\sigma_{HM}}
    \left(\frac{P^E_t}{\xi}\right)^{\sigma_{HM}} E_t.
    \label{eq:research-conditional-demands}
\end{equation}
The developer can consequently choose the scale of research in two stages. For
$\eta\in(0,1)$, its optimal scale satisfies
\begin{equation}
    P^E_t=q_t\chi\eta A_t^\phi E_t^{\eta-1},
    \qquad
    E_t=
    \left(
        \frac{q_t\chi\eta A_t^\phi}{P^E_t}
    \right)^{1/(1-\eta)}.
    \label{eq:research-scale}
\end{equation}

Define the contribution of automated services to effective research as
\begin{equation}
    s^M_t
    \equiv
    \frac{\omega_M M_t^{(\sigma_{HM}-1)/\sigma_{HM}}}
    {\omega_H H_t^{(\sigma_{HM}-1)/\sigma_{HM}}
        +\omega_M M_t^{(\sigma_{HM}-1)/\sigma_{HM}}}.
    \label{eq:automated-research-share}
\end{equation}
At an interior cost-minimizing allocation, this is also the share of research
expenditure devoted to automated services:
\begin{equation}
    s^M_t=\frac{\xi M_t}{w_tH_t+\xi M_t}.
    \label{eq:automated-research-cost-share}
\end{equation}
Equations \eqref{eq:research-input-ratio} and
\eqref{eq:automated-research-share} imply
\begin{equation}
    \frac{s^M_t}{1-s^M_t}
    =
    \left(\frac{\omega_M}{\omega_H}\right)^{\sigma_{HM}}
    \left(\frac{w_t}{\xi}\right)^{\sigma_{HM}-1}.
    \label{eq:automation-odds}
\end{equation}
An initially human-dominated research sector is an initial-condition assumption,
not a consequence of the CES specification. Define
\begin{equation}
    \Omega_0
    \equiv
    \left(\frac{\omega_M}{\omega_H}\right)^{\sigma_{HM}}
    \left(\frac{w_0}{\xi_0}\right)^{\sigma_{HM}-1}.
    \label{eq:initial-automation-odds}
\end{equation}
Then $s^M_0=\Omega_0/(1+\Omega_0)$, so $s^M_0\simeq0$ requires
\begin{equation}
    \frac{w_0}{\xi_0}
    \ll
    \left(\frac{\omega_H}{\omega_M}\right)^{\sigma_{HM}/(\sigma_{HM}-1)}.
    \label{eq:initial-human-dominance-condition}
\end{equation}
The right-hand side is the relative-price threshold at which automated services
account for one half of effective research. A low initial value of $w_0/\xi_0$
means that human researchers are inexpensive relative to a standardized unit of
automated research. A smaller CES weight $\omega_M$ reinforces this initial human
dominance. The statement depends on the units used to measure $M$ and on the
corresponding normalization of $\xi$; its economic content is a restriction on
relative unit costs, not on either nominal price in isolation.

More generally, if the unit cost of automated research varies at rate $g_{\xi,t}$,
the automation share evolves according to
\begin{equation}
    \dot s^M_t
    =
    s^M_t(1-s^M_t)(\sigma_{HM}-1)
    \left(g_{w,t}-g_{\xi,t}\right).
    \label{eq:automation-share-dynamics}
\end{equation}
The baseline holds $\xi$ constant, so $g_{\xi,t}=0$.
Integrating \eqref{eq:automation-share-dynamics} gives
\begin{equation}
    \frac{s^M_t}{1-s^M_t}
    =
    \frac{s^M_0}{1-s^M_0}
    \exp\left\{
        (\sigma_{HM}-1)\int_0^t
        \left(g_{w,\tau}-g_{\xi,\tau}\right)\dd\tau
    \right\}.
    \label{eq:integrated-automation-odds}
\end{equation}
If the relative-cost growth gap is constant at $\Delta_g>0$, the transition is
logistic:
\begin{equation}
    s^M_t
    =
    \left[
        1+\frac{1-s^M_0}{s^M_0}
        e^{-(\sigma_{HM}-1)\Delta_g t}
    \right]^{-1}.
    \label{eq:logistic-automation-transition}
\end{equation}
The date at which the automated share reaches one half is
\begin{equation}
    t_{1/2}
    =
    \frac{\log[(1-s^M_0)/s^M_0]}
    {(\sigma_{HM}-1)\Delta_g}.
    \label{eq:automation-half-time}
\end{equation}
Starting close to zero can therefore produce a long human-dominated phase even when
$s^M_t\to1$. A larger $\sigma_{HM}$ or a faster increase in the research wage relative
to automated-research cost compresses that transition.

\begin{conditionalresult}[Research mix when relative costs diverge]
\label{cond:gradual-automation}
Suppose $\sigma_{HM}>1$. If $w_t/\xi\to\infty$, then $s^M_t\to1$ and
$H_t/E_t\to0$. If $w_t/\xi$ remains bounded, $\sigma_{HM}>1$ alone does not imply that
the human contribution disappears.
\end{conditionalresult}

Conditional result~\ref{cond:gradual-automation} describes the research mix for a
given limiting relative-cost path; it does not establish that the equilibrium wage
diverges relative to $\xi$. It replaces a discrete automation threshold
with a continuous equilibrium transition. The research elasticity determines
whether human research can become asymptotically inessential. The production
elasticity $\sigma_{XL}$ affects whether that limit is reached by determining the wage
path.

\subsubsection{Automation intensity and human research employment}

The research-input ratio is not the same object as employment in the research
sector. Define the human-research share of population as $h_t\equiv H_t/N_t$.
The following identity separates research composition from research scale:
\begin{equation}
    h_t
    =
    \frac{H_t}{M_t}\frac{M_t}{N_t}.
    \label{eq:human-research-scale-identity}
\end{equation}
Equation~\eqref{eq:research-input-ratio} and population growth imply
\begin{equation}
    g_{h,t}
    =
    \sigma_{HM}\left(g_{\xi,t}-g_{w,t}\right)
    +g_{M,t}-n,
    \qquad
    g_{h,t}\equiv\frac{\dot h_t}{h_t}.
    \label{eq:human-research-share-growth}
\end{equation}

\begin{proposition}[Research automation and human research employment identities]
\label{prop:research-employment-scale}
At every interior research-input allocation,
\begin{equation}
    g_{H/M,t}
    =
    \sigma_{HM}(g_{\xi,t}-g_{w,t}),
    \qquad
    g_{H/N,t}
    =
    \sigma_{HM}(g_{\xi,t}-g_{w,t})+g_{M,t}-n.
    \label{eq:research-composition-and-scale-identities}
\end{equation}
Consequently, $H_t/M_t$ falls if and only if $g_{w,t}>g_{\xi,t}$, while the
human-research share of population rises if and only if
\begin{equation}
    g_{M,t}-n
    >
    \sigma_{HM}\left(g_{w,t}-g_{\xi,t}\right).
    \label{eq:human-research-expansion-condition}
\end{equation}
\end{proposition}

\begin{proposition}[Automated research scale, singularity, and the real wage]
\label{prop:automated-research-wage}
Suppose the research-input allocation is interior and the unit resource cost
$\xi$ is constant. At every date,
\begin{equation}
    w_t
    =
    \xi\frac{\omega_H}{\omega_M}
    \left(\frac{M_t}{H_t}\right)^{1/\sigma_{HM}}
    \geq
    \xi\frac{\omega_H}{\omega_M}
    \left(\frac{M_t}{N_t}\right)^{1/\sigma_{HM}}.
    \label{eq:wage-automated-research-lower-bound}
\end{equation}
Consequently, if automated research per capita diverges as $t\uparrow T$, where
$T$ may be finite or infinite, then $w_t\to\infty$. Any earlier finite decline in
the wage is therefore eventually reversed in levels. Wage growth also satisfies
\begin{equation}
    g_{w,t}=\frac{g_{M,t}-g_{H,t}}{\sigma_{HM}}.
    \label{eq:wage-research-scale-growth}
\end{equation}
Thus the wage is eventually increasing if automated research eventually grows
faster than human research. In addition, if $\phi<1$ and capability diverges at a
finite date $T^*$, then
\begin{equation}
    \limsup_{t\uparrow T^*}w_t=\infty.
    \label{eq:finite-singularity-wage-unbounded}
\end{equation}
Thus a finite-time capability singularity cannot occur while the wage remains
bounded.
\end{proposition}

The proposition does not infer wage recovery merely from rapid capital growth.
Equation~\eqref{eq:dynamic-wage-decomposition} shows that sufficiently rapid growth
of AI services can offset capital deepening during part of the transition. Instead,
Proposition~\ref{prop:automated-research-wage} identifies the scale of automated
research relative to the population as a sufficient statistic for divergence of
the wage. For a finite-time singularity, it also rules out a permanently bounded
wage without imposing an asymptotic resource share. These results establish
recovery in levels, not eventual monotonicity, unless $g_M>g_H$ also holds after
some date.

Proposition~\ref{prop:research-employment-scale} shows that the development industry
can become more automated while employing a larger fraction of the population. The
first force reduces the number of human researchers required per unit of automated
research. The second expands the scale of the research sector. Either force can
dominate during a transition, so $H_t/N_t$ need not be monotone even when $H_t/M_t$
falls monotonically.

Under the proportional allocation rule $\xi_tM_t=m_RY_t$, with constant $m_R$,
equation~\eqref{eq:human-research-share-growth} becomes
\begin{equation}
    g_{h,t}
    =
    g_{Y,t}-n-\sigma_{HM} g_{w,t}+(\sigma_{HM}-1)g_{\xi,t}.
    \label{eq:numerical-human-research-share-growth}
\end{equation}
The baseline holds $\xi$ constant. In that case, $H_t/N_t$ rises exactly when
per-capita output grows faster than the wage multiplied by the research-substitution
elasticity. This last condition belongs to a reduced allocation closure rather than
to the fully optimizing monopoly. It is useful for isolating a mechanism but is not
used to construct the equilibrium simulations in Section~\ref{sec:numerical}.

The private shadow value continues to satisfy
\begin{equation}
    (r_t-\delta)q_t
    =
    \pi^X_{A,t}+q_tF_{A,t}+\dot q_t,
    \qquad
    F_{A,t}=\phi\frac{\dot A_t}{A_t}.
    \label{eq:private-costate}
\end{equation}

\subsubsection{Human-dominated research}

When $s^M_t$ is close to zero,
\begin{equation}
    E_t
    \simeq
    \omega_H^{\sigma_{HM}/(\sigma_{HM}-1)}H_t.
    \label{eq:human-dominated-research}
\end{equation}
If a constant share $s_H\in(0,1)$ of the population works in research,
\begin{equation}
    H_t=s_HN_t,
    \qquad
    L_t=(1-s_H)N_t,
    \label{eq:research-share}
\end{equation}
the ideas equation implies
\begin{equation}
    g_A^H
    =
    \frac{\eta n}{1-\phi}.
    \label{eq:human-dominated-ga}
\end{equation}
This is a semi-endogenous benchmark for the early, human-dominated phase. It is not
generally a permanent balanced-growth path once automated research is feasible. If
the wage grows relative to $\xi$, equation \eqref{eq:automation-share-dynamics}
moves the economy away from this phase.

\subsubsection{AI-dominated research dynamics}

When $s^M_t\to1$, the research aggregator satisfies
\begin{equation}
    E_t
    \simeq
    \omega_M^{\sigma_{HM}/(\sigma_{HM}-1)}M_t.
    \label{eq:ai-dominated-research}
\end{equation}
The asymptotic law of motion is therefore
\begin{equation}
    \dot A_t
    =
    \kappa A_t^\phi M_t^\eta,
    \qquad
    \kappa\equiv
    \chi\omega_M^{\eta\sigma_{HM}/(\sigma_{HM}-1)}.
    \label{eq:ai-dominated-ideas}
\end{equation}
Its growth rate obeys
\begin{equation}
    \frac{\dot g_{A,t}}{g_{A,t}}
    =
    (\phi-1)g_{A,t}+\eta g_{M,t}.
    \label{eq:ai-growth-rate-dynamics}
\end{equation}
If $g_{M,t}$ converges to a positive constant and $\phi<1$, a constant capability
growth rate must satisfy
\begin{equation}
    g_A^\infty
    =
    \frac{\eta g_M^\infty}{1-\phi}.
    \label{eq:ai-dominated-ga}
\end{equation}
Research growth is no longer tied directly to population growth. It is tied to the
growth of automated-research services, whose equilibrium path depends on output,
profits, and the private value of capability.

\subsubsection{A closed-form AI-dominated path under Cobb--Douglas production}

The Cobb--Douglas limit of the production CES provides a closed-form benchmark.
Define
\begin{equation}
    \beta\equiv(1-\alpha)\omega_X,
    \qquad
    \upsilon\equiv\frac{\beta}{1-\alpha-\beta}.
    \label{eq:beta-definition}
\end{equation}
Production becomes
\begin{equation}
    Y_t=K_t^\alpha L_t^{1-\alpha-\beta}X_t^\beta.
    \label{eq:cd-production}
\end{equation}
The monopolist chooses
\begin{align}
    p^X_t&=\frac{\xi}{\beta A_t},
    &
    U_t&=\frac{\beta^2}{\xi}Y_t,
    &
    X_t&=\frac{\beta^2A_t}{\xi}Y_t,
    \label{eq:cd-monopoly}\\
    \pi^X_t&=\beta(1-\beta)Y_t.
    \label{eq:cd-monopoly-rents}
\end{align}
Substitution into production gives
\begin{equation}
    Y_t
    =
    \left(\frac{\beta^2}{\xi}\right)^{\beta/(1-\beta)}
    A_t^{\beta/(1-\beta)}
    K_t^{\alpha/(1-\beta)}
    L_t^{(1-\alpha-\beta)/(1-\beta)}.
    \label{eq:reduced-production}
\end{equation}

Let $\gamma^\infty$ denote the common growth rate of output, capital,
consumption, and inference compute per capita. Along an AI-dominated path with a
constant research-resource share $m_R\equiv\xi M_t/Y_t$,
\begin{equation}
    \gamma^\infty=\upsilon g_A^\infty,
    \qquad
    g_M^\infty=n+\gamma^\infty.
    \label{eq:cd-ai-growth-relations}
\end{equation}
Combining \eqref{eq:ai-dominated-ga} and
\eqref{eq:cd-ai-growth-relations} yields
\begin{align}
    g_A^\infty
    &=
    \frac{\eta n}
    {1-\phi-\eta\upsilon},
    \label{eq:cd-ai-ga}\\
    \gamma^\infty
    &=
    \frac{\upsilon\eta n}
    {1-\phi-\eta\upsilon}.
    \label{eq:cd-ai-gamma}
\end{align}
A positive constant-growth path requires
\begin{equation}
    1-\phi-\eta\upsilon>0.
    \label{eq:cd-ai-bg-condition}
\end{equation}
If this condition fails, the feedback from output to automated research and from
research back to output is too strong for a positive constant-growth solution. The
failure of \eqref{eq:cd-ai-bg-condition} identifies a candidate acceleration regime;
the transition system determines whether research and output accelerate or instead
move to a boundary allocation.

The remaining balanced-growth objects are
\begin{align}
    g_Y=g_K=g_C=g_U=g_M&=n+\gamma^\infty,
    \label{eq:aggregate-bg-rates}\\
    g_X&=n+\gamma^\infty+g_A^\infty,
    &
    g_w&=\gamma^\infty,
    &
    g_{p^X}&=-g_A^\infty,
    \label{eq:other-bg-rates}\\
    r^\infty&=\rho+\delta+\gamma^\infty,
    &
    \frac{K_t}{Y_t}
    &=
    \frac{\alpha}{\rho+\delta+\gamma^\infty}.
    \label{eq:bg-capital-output}
\end{align}
The costate equation and the research first-order condition imply
\begin{equation}
    m_R^\infty
    =
    \frac{\beta^2\eta g_A^\infty}
    {\rho-n+(1-\phi)g_A^\infty}.
    \label{eq:bg-research-resource-share}
\end{equation}
Consequently,
\begin{align}
    \frac{I_t}{Y_t}
    &=
    \frac{\alpha(n+\gamma^\infty+\delta)}
    {\rho+\delta+\gamma^\infty},
    \label{eq:bg-investment-share}\\
    \frac{C_t}{Y_t}
    &=
    1-
    \frac{\alpha(n+\gamma^\infty+\delta)}
    {\rho+\delta+\gamma^\infty}
    -\beta^2-m_R^\infty.
    \label{eq:bg-consumption-share}
\end{align}
An interior path requires a positive consumption share.

\begin{conditionalresult}[Candidate AI-dominated limit under Cobb--Douglas production]
\label{cond:ai-dominated-cd-limit}
Suppose $\sigma_{XL}=1$, $n>0$, $\xi$ is constant, and consider a nondegenerate
constant-growth limit in which $s^M_t\to1$, $g_A>0$, $\xi M/Y$ is constant, and the
household and developer optimality conditions hold. Then the capability and
per-capita growth rates are given by \eqref{eq:cd-ai-ga} and
\eqref{eq:cd-ai-gamma}; the remaining growth rates, capital--output ratio, research
share, investment share, and consumption share are given by
\eqref{eq:aggregate-bg-rates}--\eqref{eq:bg-consumption-share}. Such a limit requires
\begin{equation}
    D_{\mathrm{AI}}
    =1-\phi-\eta\frac{\omega_X}{\omega_L}>0
    \label{eq:ai-dominated-cd-necessary-stability}
\end{equation}
and the consumption share in \eqref{eq:bg-consumption-share} to be positive. If
$D_{\mathrm{AI}}\leq0$, no AI-dominated constant-growth limit with a positive finite
$g_A$ exists.
\end{conditionalresult}

This is a necessary characterization of the stated candidate limit. It does
not establish that a global equilibrium path exists or converges to that limit; those
claims additionally require the transition system and transversality conditions.

Because $g_w=\gamma^\infty>0$ and $\xi$ is constant, conditional
result~\ref{cond:gradual-automation} applies to this candidate. More precisely,
\begin{equation}
    g_{H/N}
    =
    -(\sigma_{HM}-1)\gamma^\infty<0.
    \label{eq:human-research-population-share}
\end{equation}
The fraction of the population employed in frontier research therefore converges to
zero. The number of human researchers need not fall in levels: it grows at rate
$n-(\sigma_{HM}-1)\gamma^\infty$ and declines only if the second term exceeds population
growth.

\subsubsection{Human essentiality and the growth regime}

The CES aggregator provides a unified way to state the feedback condition during
the transition. Because it is constant returns to scale, its exact growth
decomposition is
\begin{equation}
    g_{E,t}
    =
    (1-s^M_t)g_{H,t}+s^M_tg_{M,t}.
    \label{eq:ces-research-growth-decomposition}
\end{equation}
Consider a reduced proportional-allocation closure with $g_H=n$,
$g_M=n+\gamma_t$, and $\gamma_t=\upsilon g_{A,t}$. These relations hold on the
Cobb--Douglas balanced-growth path and are extended here to the transition only to
isolate the feedback mechanism. Equations \eqref{eq:ideas} and
\eqref{eq:ces-research-growth-decomposition} imply
\begin{align}
    \dot g_{A,t}
    &=
    g_{A,t}\left[\eta n-D(s^M_t)g_{A,t}\right],
    \label{eq:ces-transition-growth-dynamics}\\
    D(s)
    &\equiv
    1-\phi-\eta\upsilon s.
    \label{eq:automation-feedback-denominator}
\end{align}
Human-dominated research has $D(0)=1-\phi$. The AI-dominated limit has
$D(1)=1-\phi-\eta\upsilon$. Research automation therefore reduces the force that
stabilizes capability growth.
Under the same closure, cost minimization makes the relative-price path consistent
with the input-growth rates:
\begin{equation}
    g_{w,t}-g_{\xi,t}
    =
    \frac{g_{M,t}-g_{H,t}}{\sigma_{HM}}
    =
    \frac{\upsilon}{\sigma_{HM}}g_{A,t}.
    \label{eq:consistent-relative-cost-growth}
\end{equation}
The automation share consequently satisfies
\begin{equation}
    \dot s^M_t
    =
    s^M_t(1-s^M_t)
    \frac{\sigma_{HM}-1}{\sigma_{HM}}\upsilon g_{A,t}.
    \label{eq:reduced-consistent-automation-dynamics}
\end{equation}
Equations \eqref{eq:ces-transition-growth-dynamics} and
\eqref{eq:reduced-consistent-automation-dynamics} form a two-dimensional reduced
system for the automation share and capability growth. They do not impose an
independent relative-wage path.

\subsubsection{The fully Cobb--Douglas particular case}

The AI-dominated path is not the Cobb--Douglas limit of the research technology.
To isolate the case in which human research remains essential, set
$\sigma_{XL}=\sigma_{HM}=1$. The two technologies become
\begin{align}
    Y_t
    &=K_t^\alpha L_t^{(1-\alpha)\omega_L}
        X_t^{(1-\alpha)\omega_X},
    \label{eq:fully-cd-production}\\
    E_t^{\mathrm{CD}}
    &=H_t^{\omega_H}M_t^{\omega_M}.
    \label{eq:cd-research-counterfactual}
\end{align}
Both research inputs are essential. The monopoly solution can be written directly
in terms of the primitive production weights:
\begin{equation}
    U_t
    =
    \frac{[(1-\alpha)\omega_X]^2}{\xi}Y_t,
    \qquad
    X_t
    =
    \frac{[(1-\alpha)\omega_X]^2}{\xi}A_tY_t.
    \label{eq:fully-cd-monopoly-primitives}
\end{equation}

Consider an interior constant-growth path. Constant population allocations imply
$g_H=g_L=n$, while constant automated-research expenditure relative to output implies
$g_M=g_Y$. Production and research growth then satisfy
\begin{align}
    g_Y-n
    &=
    \frac{\omega_X}{\omega_L}g_A,
    \label{eq:fully-cd-output-capability-growth}\\
    g_E^{\mathrm{CD}}
    &=
    n+\omega_M(g_Y-n).
    \label{eq:cd-research-growth}
\end{align}
The first equation is the production feedback: capability raises services per unit
of compute, services expand output, and the output expansion scales the resources
available to automated research. The second equation shows that only the
$\omega_M$ contribution of faster-growing automated research enters effective
research when the human input remains essential.

Log-differentiating the ideas equation and imposing constant $g_A$ gives
$(1-\phi)g_A=\eta g_E$. Combining this condition with
\eqref{eq:fully-cd-output-capability-growth}--\eqref{eq:cd-research-growth} yields
\begin{align}
    g_A^{\mathrm{CD}}
    &=
    \frac{\eta n}
    {1-\phi-\eta\omega_M\frac{\omega_X}{\omega_L}},
    \label{eq:full-cd-ga}\\
    g_Y^{\mathrm{CD}}-n
    &=
    \frac{\omega_X}{\omega_L}
    \frac{\eta n}
    {1-\phi-\eta\omega_M\frac{\omega_X}{\omega_L}}.
    \label{eq:full-cd-gamma}
\end{align}
The remaining growth rates are
\begin{align}
    g_H=g_L=g_N&=n,
    \notag\\
    g_K=g_C=g_U=g_M
    &=g_Y^{\mathrm{CD}},
    \notag\\
    g_X&=g_Y^{\mathrm{CD}}+g_A^{\mathrm{CD}},
    &
    g_w&=g_Y^{\mathrm{CD}}-n,
    &
    g_{p^X}&=-g_A^{\mathrm{CD}}.
    \label{eq:fully-cd-growth-rates}
\end{align}
Thus $H/N$ and $L/N$ are constant, but $H/M$ falls at rate
$n-g_Y^{\mathrm{CD}}<0$. Human research becomes small relative to automated research
without ceasing to be essential.

The denominator in \eqref{eq:full-cd-ga}--\eqref{eq:full-cd-gamma} is positive by
Assumption~\ref{ass:stable-human-essential-benchmark}. The fully Cobb--Douglas
candidate therefore has positive finite growth rates.
The term $1-\phi$ measures the force of diminishing returns to the existing
capability stock. The term
$\eta\omega_M\omega_X/\omega_L$ measures the complete feedback loop. A one-percent
increase in capability raises per-capita output by $\omega_X/\omega_L$ percent;
the resulting expansion of automated research affects effective research with
weight $\omega_M$ and affects capability growth with elasticity $\eta$. The maintained
restriction selects the parameter region in which diminishing returns dominate this
loop in the baseline. Conditional result~\ref{cond:singularity-conditions}
deliberately relaxes the restriction to show what occurs at and beyond its boundary.

The condition in \eqref{eq:fully-cd-stability-condition} is not sufficient for an
interior allocation. The Euler equation, capital accumulation, and the developer's
research condition imply
\begin{align}
    \frac{K_t}{Y_t}
    &=
    \frac{\alpha}{\rho+\delta+g_Y^{\mathrm{CD}}-n},
    \label{eq:fully-cd-capital-output}\\
    \frac{I_t}{Y_t}
    &=
    \frac{\alpha(g_Y^{\mathrm{CD}}+\delta)}
    {\rho+\delta+g_Y^{\mathrm{CD}}-n},
    \label{eq:fully-cd-investment-share}\\
    \frac{\xi U_t}{Y_t}
    &=
    [(1-\alpha)\omega_X]^2,
    \label{eq:fully-cd-inference-share}\\
    \frac{\xi M_t}{Y_t}
    &=
    \frac{[(1-\alpha)\omega_X]^2\eta\omega_M g_A^{\mathrm{CD}}}
    {\rho-n+(1-\phi)g_A^{\mathrm{CD}}}.
    \label{eq:full-cd-research-resource-share}
\end{align}
The resource constraint requires $C_t/Y_t>0$. Substituting the growth rates and
resource shares above gives the following restriction entirely in primitive
parameters:
\begin{equation}
\begin{aligned}
&\frac{\alpha\left\{(n+\delta)
    \left(1-\phi-\eta\omega_M\frac{\omega_X}{\omega_L}\right)
    +\eta n\frac{\omega_X}{\omega_L}\right\}}
{(\rho+\delta)
    \left(1-\phi-\eta\omega_M\frac{\omega_X}{\omega_L}\right)
    +\eta n\frac{\omega_X}{\omega_L}}
+[(1-\alpha)\omega_X]^2
\\
&\qquad
+\frac{[(1-\alpha)\omega_X]^2\eta^2\omega_M n}
{(\rho-n)
    \left(1-\phi-\eta\omega_M\frac{\omega_X}{\omega_L}\right)
    +(1-\phi)\eta n}
<1.
\end{aligned}
\label{eq:fully-cd-consumption-feasibility}
\end{equation}
The three terms on the left are, respectively, the output shares absorbed by
physical investment, inference compute, and automated research. The inequality is
equivalent to a positive consumption share when
\eqref{eq:fully-cd-stability-condition} holds. The allocation of workers is also
constant because
\begin{equation}
    \frac{H_t}{L_t}
    =
    \frac{\omega_H}{\omega_M}
    \frac{\xi M_t/Y_t}{(1-\alpha)\omega_L}.
    \label{eq:fully-cd-labor-allocation}
\end{equation}

For comparison with the AI-dominated research limit, define
\begin{equation}
    D_{\mathrm{CD}}
    \equiv
    1-\phi-\eta\omega_M\frac{\omega_X}{\omega_L},
    \qquad
    D_{\mathrm{AI}}
    \equiv
    1-\phi-\eta\frac{\omega_X}{\omega_L}.
    \label{eq:growth-feedback-denominators}
\end{equation}

\begin{conditionalresult}[Comparison of the two Cobb--Douglas candidate limits]
\label{cond:human-essentiality-acceleration}
Consider the two candidate constant-growth allocations derived above. The fully
Cobb--Douglas candidate maintains constant population allocations and a constant
automated-research expenditure share. The AI-dominated candidate satisfies the
conditions in conditional result~\ref{cond:ai-dominated-cd-limit}. The first candidate
requires \eqref{eq:fully-cd-consumption-feasibility}; the second requires
$D_{\mathrm{AI}}>0$. Their feedback denominators satisfy
\begin{equation}
    D_{\mathrm{CD}}-D_{\mathrm{AI}}
    =
    \eta\omega_H\frac{\omega_X}{\omega_L}>0.
    \label{eq:feedback-denominator-gap}
\end{equation}
There is consequently a nonempty parameter region with
$D_{\mathrm{CD}}>0\geq D_{\mathrm{AI}}$: the fully Cobb--Douglas candidate has
positive finite growth rates even though no positive constant-growth allocation
exists within the stated AI-dominated candidate class.
\end{conditionalresult}

For Cobb--Douglas research, the automated-research contribution is fixed at $\omega_M$,
so \eqref{eq:ces-transition-growth-dynamics} becomes
\begin{equation}
    \dot g_{A,t}
    =
    g_{A,t}\left(\eta n-D_{\mathrm{CD}}g_{A,t}\right).
    \label{eq:full-cd-growth-dynamics}
\end{equation}
Assumption~\ref{ass:stable-human-essential-benchmark} implies that
\eqref{eq:full-cd-growth-dynamics} has the stable positive fixed point
$\eta n/D_{\mathrm{CD}}$ in the baseline. The following result temporarily relaxes
the assumption. It first shows that the fully Cobb--Douglas economy can accelerate
when $D_{\mathrm{CD}}\leq0$, even though $\sigma_{XL}=1$. It then returns to the
baseline CES research technology, whose denominator can become nonpositive as
research is automated even when $D_{\mathrm{CD}}>0$.

\begin{definition}[Finite-time singularity]
Following \citet{aghionjonesjones2019}, we reserve the term singularity for an
unbounded outcome reached in finite time. In this model, a finite-time singularity
occurs if there is a finite date $T^*$ such that $g_{A,t}$ and $A_t$ diverge as
$t\uparrow T^*$. Accelerating growth that remains finite at every finite date is not
a singularity.
\end{definition}

\begin{conditionalresult}[Recursive-feedback regimes at unit production elasticity]
\label{cond:singularity-conditions}
Suppose $g_{A,0}>0$, $n>0$, $\sigma_{XL}=1$, and the reduced
proportional-allocation closure in \eqref{eq:ces-transition-growth-dynamics} holds.
When $\sigma_{HM}=1$, the automated-research contribution is the constant
$\omega_M$ and $D(s^M_t)=D_{\mathrm{CD}}$. Then:
\begin{enumerate}[label=(\roman*)]
    \item If $D_{\mathrm{CD}}>0$, $g_A$ converges to the stable rate
    $\eta n/D_{\mathrm{CD}}$.
    \item If $D_{\mathrm{CD}}=0$, $g_A$ grows exponentially and $A$ grows double
    exponentially, but both remain finite at every finite date.
    \item If $D_{\mathrm{CD}}<0$, the reduced economy reaches a finite-time
    singularity at
    \begin{equation}
        T^*
        =
        \frac{1}{\eta n}
        \log\left[
            1+\frac{\eta n}{(-D_{\mathrm{CD}})g_{A,0}}
        \right].
        \label{eq:constant-feedback-singularity-time}
    \end{equation}
    \pagebreak
    \item If instead $\sigma_{HM}>1$, the contribution $s^M_t$ is endogenous. No
    finite-time singularity occurs if $D(s^M_t)\geq0$ for every $t$. A singularity
    occurs if and only if the bracket in
    \begin{equation}
        \frac{1}{g_{A,t}}
        =
        e^{-\eta nt}
        \left[
            \frac{1}{g_{A,0}}
            +\int_0^t e^{\eta n\tau}D(s^M_\tau)\dd\tau
        \right]
        \label{eq:time-varying-singularity-condition}
    \end{equation}
    reaches zero at a finite date.
\end{enumerate}
\end{conditionalresult}

Parts (ii) and (iii) establish that explosive growth does not require
$\sigma_{XL}>1$. At unit substitution elasticities, sufficiently strong recursive
feedback is enough. Part (ii) produces an unbounded growth rate only as
$t\to\infty$; part (iii) produces a finite-time singularity. The relevant threshold
is zero, not one: every $D_{\mathrm{CD}}>0$ stabilizes the reduced system, although
a small positive value implies a high limiting growth rate.

The reduced-system calculation separates three objects that are sometimes conflated. A positive
constant-growth path requires $D>0$. The knife-edge $D=0$ eliminates balanced
growth but does not create a finite-time singularity. A negative and persistent
$D$ does. In the CES transition, the instantaneous threshold is
\begin{equation}
    s^\dagger
    =
    \frac{1-\phi}{\eta\upsilon}.
    \label{eq:critical-automation-share}
\end{equation}
When $D_{\mathrm{AI}}>0$, $D(s)>0$ for every $s\in[0,1]$ and the reduced system
cannot produce a singularity. When $D_{\mathrm{AI}}=0$, the damping term converges
to zero but remains nonnegative, so there is no finite-time singularity. When
$D_{\mathrm{AI}}<0$ and $s^M_t\to1$, the automation share eventually exceeds
$s^\dagger$ and remains bounded above it; the integral in
\eqref{eq:time-varying-singularity-condition} then reaches zero at a finite date.
A temporary interval with $D(s^M_t)<0$ is not sufficient by itself: the exact
condition is the accumulated, exponentially weighted integral in
\eqref{eq:time-varying-singularity-condition}.

These are singularity conditions for the reduced closure, not a global explosion
theorem for the fully optimizing monopoly. The proportional spending rule, the
capital--output relation, or the interior consumption allocation may fail before
$T^*$. A full equilibrium can therefore approach a resource constraint or a
boundary allocation instead. The denominator $D_{\mathrm{AI}}<0$ should be read as
a candidate singularity region whose equilibrium relevance must be checked with the
complete transition system.

The comparison in conditional result~\ref{cond:human-essentiality-acceleration} does not
assume that the research elasticity changes over time. In the baseline CES model,
the fixed elasticity $\sigma_{HM}>1$ permits a transition from the human-dominated growth
rate in \eqref{eq:human-dominated-ga} to the AI-dominated dynamics in
\eqref{eq:ai-growth-rate-dynamics}. The Cobb--Douglas counterfactual asks what would
happen if the human input remained essential throughout. Under the baseline
calibration, $D_{\mathrm{CD}}=0.3106$ and $D_{\mathrm{AI}}=0.2375$, so both reduced
candidate growth rates are positive and finite. Holding $\eta=0.45$ and the CES
weights at their baseline values, raising $\phi$ to $0.90$ gives
$D_{\mathrm{CD}}=0.060625$ but $D_{\mathrm{AI}}=-0.0125$. The human-essential
reduced feedback equation then has a positive finite fixed point, whereas the
AI-dominated reduced dynamics enter the candidate acceleration region. This
comparison does not establish existence of an optimizing equilibrium in either
case.

\subsubsection{Why the production CES matters}

The Cobb--Douglas case is the knife-edge at which expenditure shares remain
constant. For $\sigma_{XL}\neq1$, the monopoly condition
\eqref{eq:monopoly-ai-price} implies different asymptotic regimes as the marginal
cost $\xi/A_t$ falls.

\begin{conditionalresult}[Asymptotic monopoly supply under increasingly capable AI]
\label{cond:production-ces-regimes}
Suppose $A_t\to\infty$, the service-market solution remains interior, and the
effective marginal cost of AI services vanishes relative to capital intensity:
\begin{equation}
    \frac{\xi}{A_t(K_t/L_t)^\alpha}\longrightarrow0.
    \label{eq:vanishing-normalized-ai-cost}
\end{equation}
\begin{enumerate}[label=(\roman*)]
    \item If $\sigma_{XL}<1$, the monopolist's AI share converges to
    \begin{equation}
        \bar s^X
        =
        \frac{1-\sigma_{XL}}{1-\alpha\sigma_{XL}}.
        \label{eq:low-sigma-ai-share}
    \end{equation}
    The ratio $X_t/L_t$ remains bounded and the productive effect of further
    reductions in $\xi/A_t$ vanishes.
    \item If $\sigma_{XL}=1$, factor shares are constant and any AI-dominated
    balanced-growth candidate, when it exists, is characterized by
    \eqref{eq:cd-ai-ga}--\eqref{eq:bg-consumption-share}.
    \item If $\sigma_{XL}>1$, $s^X_t\to1$, the production-labor share converges to zero,
    and
    \begin{equation}
        \frac{Y_t}{K_t}
        \sim
        D_{XL} A_t^{(1-\alpha)/\alpha},
        \qquad D_{XL}>0.
        \label{eq:high-sigma-ak-limit}
    \end{equation}
    The output--capital ratio and the return to capital therefore diverge with
    capability.
\end{enumerate}
\end{conditionalresult}

The condition in \eqref{eq:vanishing-normalized-ai-cost} is satisfied whenever
$A_t\to\infty$ and $K_t/L_t$ does not collapse too quickly. For
$\sigma_{XL}>1$, it implies $X_t/L_t\to\infty$ and therefore $s^X_t\to1$.
Appendix~\ref{sec:appendix-proofs} provides the proof and derives the constant
$D_{XL}$ in \eqref{eq:high-sigma-ak-limit}.

The apparent discontinuity at unit elasticity is real as an asymptotic statement,
but it is not a discontinuity in finite allocations. The distinction follows from
the exact CES share equation.

\begin{proposition}[The unit-elasticity singular perturbation]
\label{prop:production-unit-elasticity-knife-edge}
Let $\mathcal R_t\equiv w_t/p^X_t$ denote the relative price of production labor
and suppose $\omega_L,\omega_X>0$ with
$\omega_L+\omega_X=1$. At every interior allocation,
\begin{equation}
    \frac{s^X_t}{1-s^X_t}
    =
    \left(\frac{\omega_X}{\omega_L}\right)^{\sigma_{XL}}
    \mathcal R_t^{\sigma_{XL}-1}.
    \label{eq:production-share-odds}
\end{equation}
Let $s^X(\mathcal R,\sigma)$ denote the share defined by the right-hand side of
\eqref{eq:production-share-odds}. Then:
\begin{enumerate}[label=(\roman*)]
    \item for every finite $\mathcal R>0$, $s^X(\mathcal R,\sigma)$ is continuous in $\sigma$ at one
    and converges to $\omega_X$ as $\sigma\to1$;
    \item $s^X(\mathcal R,1)=\omega_X$ for every $\mathcal R$, whereas
    $\lim_{\mathcal R\to\infty}s^X(\mathcal R,\sigma)=1$ for every fixed $\sigma>1$; hence
    \begin{equation}
            \lim_{\sigma_{XL}\downarrow1}\lim_{\mathcal R\to\infty}s^X(\mathcal R,\sigma_{XL})=1,
        \qquad
            \lim_{\mathcal R\to\infty}\lim_{\sigma_{XL}\downarrow1}s^X(\mathcal R,\sigma_{XL})
        =\omega_X;
        \label{eq:noncommuting-production-limits}
    \end{equation}
    \item for any target $\bar s\in(\omega_X,1)$ and
    $\sigma>1$, the relative price required to attain $s^X=\bar s$ is
    \begin{equation}
        \mathcal R_{\bar s}(\sigma)
        =
        \left[
            \frac{\bar s}{1-\bar s}
            \left(\frac{\omega_L}{\omega_X}\right)^{\sigma}
        \right]^{1/(\sigma-1)},
        \label{eq:production-share-threshold}
    \end{equation}
    and $\mathcal R_{\bar s}(\sigma)\to\infty$ as $\sigma\downarrow1$.
\end{enumerate}
\end{proposition}

Proposition~\ref{prop:production-unit-elasticity-knife-edge} explains why the
Cobb--Douglas and gross-substitution asymptotics are not contradictory. At
$\sigma_{XL}=1$, the relative price affects the input ratio but not the CES
contribution share. For every fixed elasticity above one, however small the gap,
the share function converges to one as $w/p^X$ becomes unbounded. Establishing that
the equilibrium generates this relative-price limit is a separate dynamic question. The limits
in \eqref{eq:noncommuting-production-limits} do not commute. This is a singular
perturbation: equilibrium allocations are continuous in $\sigma_{XL}$ at every
finite relative price and date, while their infinite-horizon classification is not
uniform near one.

The scale of the transition becomes arbitrarily large as the elasticity approaches
one. With $\omega_X=0.2$, for example, reaching $s^X=1/2$ requires
$\log_{10}(w/p^X)=60.81$ when $\sigma_{XL}=1.01$, $12.64$ when it is $1.05$, and
$6.62$ when it is $1.10$. The AI-dominated coefficient in
\eqref{eq:high-sigma-ak-limit} also satisfies
$D_{XL}\to0$ as $\sigma_{XL}\downarrow1$ because $\omega_X\in(0,1)$. Both facts
warn against using the asymptotic regime as a medium-run approximation when the
elasticity is close to one. This calculation characterizes the analytical scale of
the asymptotic transition; it is not a simulated equilibrium path.

Assumption~\ref{ass:stable-human-essential-benchmark} stabilizes the fully
Cobb--Douglas research feedback at exactly $\sigma_{XL}=1$. It does not remove the
separate production-side mechanism for a fixed $\sigma_{XL}>1$ and unbounded
capability. The first condition concerns feedback through automated research; the
second concerns the way the production CES converts an ever lower effective AI
price into the return to capital.

The exact equilibrium identities permit a sharper nonexistence result than merely
ruling out a constant capital--output ratio. We use the paper's definition of an
interior steady state: a path on which all aggregate variables eventually grow at
finite constant rates.

\newpage
\begin{proposition}[Gross production substitution rules out positive-growth steady state]
\label{prop:gross-substitution-no-steady-state}
Suppose $\sigma_{XL}>1$ and the service-market solution and household allocation
remain interior. There is no steady state on which all aggregate variables grow at
finite constant rates and capability grows at a positive constant rate.
\end{proposition}

The result does not require $s^X_t\to1$. The exact equilibrium identity derived in
Appendix~\ref{sec:appendix-proofs} gives a positive lower bound for
$(Y_t/K_t)/A_t^{(1-\alpha)/\alpha}$ for every interior $s^X_t\in(0,1)$. The proof
then uses the household Euler equation to rule out finite constant consumption
growth. If \eqref{eq:vanishing-normalized-ai-cost} also holds, the stronger limit
$s^X_t\to1$ delivers the asymptotic equality
\eqref{eq:high-sigma-ak-limit}.

The proof assumes a positive-growth steady state only to obtain a contradiction:
positive constant $g_A$ implies $A_t\to\infty$, which makes the return to capital and,
through the Euler equation, consumption growth diverge. The inequality
$\sigma_{XL}>1$ does not exclude a boundary path on which research shuts down and
capability remains bounded. Conditional on a finite capability stock, the production
side becomes a standard constant-returns economy and does not mechanically rule out
a stationary per-capita allocation. Whether that boundary is an equilibrium depends
on the developer's research optimality conditions and the private value of capability.

The absence of a positive-growth steady state leaves open the form of an accelerating
equilibrium. The following calculation characterizes the limiting ratios of a
postulated interior path that reaches an AI-dominated finite-time singularity. It
does not prove that such a path exists from the model's initial conditions.

Define
\begin{equation}
    \vartheta\equiv\frac{1-\alpha}{\alpha},
    \qquad
    d_U\equiv(1-\alpha)^2.
    \label{eq:high-sigma-singular-definitions}
\end{equation}
The first term is the capability elasticity of $Y/K$ in the AI-dominated
production limit. The second is the limiting share of output used for inference
compute under monopoly.

\begin{conditionalresult}[Singular scaling on a postulated AI-dominated branch]
\label{cond:high-sigma-singular-scaling}
Suppose $1<\sigma_{XL}<1/\alpha$, $\sigma_{HM}>1$, and consider an interior
equilibrium that reaches a finite date $T^*$ with $A_t\to\infty$,
$s^X_t\to1$, and $s^M_t\to1$. Suppose also that $C/Y$, investment and
automated-research resource shares, $g_A/(Y/K)$, and $qA/K$ converge to positive
finite limits. In addition, suppose that these convergences are asymptotically
regular: the growth rate of each convergent ratio, divided by $Y/K$, tends to zero.
Then
\begin{align}
    \frac{g_A}{Y/K}
    &\longrightarrow
    \bar a
    \equiv
    \frac{\eta\alpha}{1-\phi+\vartheta},
    &
    \frac{qA}{K}
    &\longrightarrow
    \frac{d_U}{\eta\alpha},
    \label{eq:high-sigma-singular-a-q}\\
    \frac{\xi M}{Y}
    &\longrightarrow
    \bar m_R
    \equiv
    \frac{\eta d_U}{1-\phi+\vartheta},
    &
    \frac{I}{Y}
    &\longrightarrow
    \bar i
    \equiv
    \alpha-\vartheta\bar a,
    \label{eq:high-sigma-singular-resource-shares}\\
    \frac{C}{Y}
    &\longrightarrow
    \bar c
    \equiv
    1-d_U-\bar m_R-\bar i.
    \label{eq:high-sigma-singular-consumption-share}
\end{align}
If $\bar c>0$ and $\bar i>0$, the output--capital ratio satisfies
\begin{equation}
    \frac{Y_t}{K_t}
    \sim
    \frac{1}{\vartheta\bar a\,(T^*-t)}.
    \label{eq:high-sigma-time-to-singularity}
\end{equation}
Thus the output--capital ratio, the return to capital, and growth rates diverge at
the finite date $T^*$ along this conditional limit. At the same time,
$K/C\to0$ and $qA/C\to0$, so neither capital nor capability retains a positive
terminal value relative to consumption at the singular boundary. Output and
automated research per capita diverge, and Proposition~\ref{prop:automated-research-wage}
therefore implies
\begin{equation}
    \frac{Y_t}{N_t}\longrightarrow\infty,
    \qquad
    \frac{M_t}{N_t}\longrightarrow\infty,
    \qquad
    w_t\longrightarrow\infty.
    \label{eq:singular-wage-divergence}
\end{equation}
Hence any finite wage decline along this conditional branch is eventually reversed
in levels before $T^*$.
\end{conditionalresult}

The proof is in Appendix~\ref{sec:appendix-proofs}. The limiting ratios depend on
the capital, capability, and research elasticities, but not directly on the value
of $\sigma_{XL}$ within the stated region. The production elasticity instead has
a large effect on when the economy reaches the AI-dominated region. This distinction
motivates the free-boundary equilibrium computations in Section~\ref{sec:numerical}.

The nonexistence result can be paired with a conditional characterization of
constant growth on the other side of the Cobb--Douglas threshold.

\begin{conditionalresult}[Necessary rates of candidate non-unit-elasticity limits]
\label{cond:nonunit-ces-steady-states}
Suppose $n>0$, $A_t\to\infty$, $\xi$ is constant, and the service-market solution
remains interior.
\begin{enumerate}[label=(\roman*)]
    \item Suppose $\sigma_{XL}<1$, \eqref{eq:vanishing-normalized-ai-cost} holds,
    and consider a nondegenerate constant-growth limit of the fully optimizing
    equilibrium. In particular, $K/Y$ and $C/Y$ converge to positive constants,
    $L/N$ converges to a positive constant, capability grows at a positive constant
    rate, human and automated research remain positive at every finite date, and the
    developer's first-order and costate conditions hold. Then
    \begin{equation}
        g_X=g_L=g_Y=g_K=g_C=n,
        \qquad
        g_A=\frac{\eta n}{1-\phi+\eta},
        \label{eq:low-sigma-optimizing-growth-rates}
    \end{equation}
    and
    \begin{equation}
        g_H=g_M=g_E
        =\frac{(1-\phi)n}{1-\phi+\eta},
        \qquad
        g_U=n-g_A.
        \label{eq:low-sigma-optimizing-research-rates}
    \end{equation}
    Output, capital, and consumption per capita are constant. The real wage, the
    return to capital, $X/L$, and the research mix are also constant. The AI-services
    contribution converges to \eqref{eq:low-sigma-ai-share}, while the resource cost
    of inference, $\xi U/Y$, converges to zero. Human and automated research grow
    more slowly than population and output:
    \begin{equation}
        g_{H/N}=g_{\xi M/Y}=-g_A<0.
        \label{eq:low-sigma-declining-research-shares}
    \end{equation}
    \item If $\sigma_{XL}>1$ and consumption is positive, there is no interior
    steady state with unbounded capability on which all aggregate variables grow at
    finite constant rates.
\end{enumerate}
\end{conditionalresult}

Combining this result with the two unit-elasticity cases gives a concise wage
classification for the constant-growth limits characterized in the paper. Under
$\sigma_{XL}<1$, the real wage is constant. Under $\sigma_{XL}=1$, it grows at the
positive rate of output per capita whenever capability improves along either the
human-essential Cobb--Douglas candidate or the AI-dominated candidate. Under
$\sigma_{XL}>1$, Proposition~\ref{prop:gross-substitution-no-steady-state} rules out
an interior steady state with improving capability and finite constant growth
rates. Thus none of the constant-growth limits characterized here has a declining
real wage; the statement is weak rather than strict because the complementarity
case has zero wage growth.

The proof is in Appendix~\ref{sec:appendix-proofs}. The first result differs from a
fixed-research-share mechanism closure. Holding $H/N$ and
$\xi M/Y$ constant would instead give $g_E=n$ and
$g_A=\eta n/(1-\phi)$. That technological path does not satisfy the developer's
intertemporal optimality conditions when $\sigma_{XL}<1$: as the productive value of
additional capability vanishes, an optimizing developer lets research grow more
slowly than output. Equations~\eqref{eq:low-sigma-optimizing-growth-rates}--
\eqref{eq:low-sigma-declining-research-shares} are necessary conditions, not a
general existence theorem. The simulations in Section~\ref{sec:numerical} instead
impose the developer's first-order and costate conditions. Existence still requires
the level first-order conditions, transversality, a feasible labor allocation, and
$C/Y>0$.

When $\sigma_{XL}=1$, the positive wage growth in
\eqref{eq:other-bg-rates} makes automated research asymptotically dominant if
$\sigma_{HM}>1$. When $\sigma_{XL}>1$, production itself becomes an acceleration
mechanism. Research automation reinforces that mechanism, but the wage path must be
solved jointly with capital accumulation before concluding that the human research
share disappears.

The stabilization comparison in conditional
result~\ref{cond:human-essentiality-acceleration}
is therefore conditional on Cobb--Douglas production. Recall that
$\vartheta=(1-\alpha)/\alpha>0$ from
\eqref{eq:high-sigma-singular-definitions}.
If research is Cobb--Douglas, $E_t=H_t^{\omega_H}M_t^{\omega_M}$, the automated
contribution share remains equal to $\omega_M$ and the human input remains essential.
Nevertheless,
equation~\eqref{eq:high-sigma-ak-limit} implies, under a fixed positive investment
share $s_I$, that
\begin{equation}
    g_{K,t}
    =s_I\frac{Y_t}{K_t}-\delta
    \sim s_ID_{XL} A_t^\vartheta-\delta.
    \label{eq:sigma-hm-one-high-sigma-xl-capital-growth}
\end{equation}
Thus, if $A_t\to\infty$, capital growth cannot converge to a finite constant.

\begin{conditionalresult}[Fixed-share mechanism with human-essential research]
\label{cond:sigma-hm-one-high-sigma-xl}
Suppose $\sigma_{HM}=1$, $\sigma_{XL}>1$, $A_t\to\infty$, and the interior monopoly solution
continues to apply. With a fixed positive investment share, there is no balanced-growth
path on which all aggregate variables grow at finite constant rates. Under the
additional reduced-form closure
\begin{equation}
    H_t=hN_t,\qquad
    \xi M_t=m_RY_t,
    \qquad h,m_R>0,
    \label{eq:sigma-hm-one-fixed-share-closure}
\end{equation}
and constant $\xi$, any $\omega_M>0$ preserves a feedback from the scale of output to
capability. The corresponding high-capability asymptotic system reaches a finite-time
singularity.
\end{conditionalresult}

This fixed-share calculation separates two meanings of human essentiality. At $\sigma_{HM}=1$, the
human contribution to each unit of effective research never disappears. But a
positive automated-input elasticity $\omega_M$ still lets the expanding scale of output
raise $M_t$ and therefore research productivity. If $\omega_M=0$, production growth still
accelerates when $\sigma_{XL}>1$, but the finite-time research feedback in the reduced
system disappears. The singularity result is conditional on fixed positive saving
and research shares. It is not yet a theorem for the fully optimizing equilibrium:
endogenous expenditure shares, a boundary at $H_t=N_t$, rising compute costs, or
physical constraints may interrupt the feedback.

\subsubsection{Private and social research}

At the social allocation, the planner's first-order condition for research compute is
\begin{equation}
    Q_tF_{M,t}=\xi.
    \label{eq:social-compute-foc}
\end{equation}
The private condition is $q_tF_{M,t}=\xi$. If $Q_t$ denotes the marginal social
value of capability evaluated consistently at the private allocation, the marginal
social net benefit of additional research compute there is
\begin{equation}
    Q_tF_{M,t}-\xi=(Q_t-q_t)F_{M,t}.
    \label{eq:knowledge-wedge}
\end{equation}

\begin{conditionalresult}[Shadow-value comparison at the private allocation]
\label{cond:knowledge-wedge}
If $Q_t>q_t$, the integrated monopolist invests too little in research compute at the
margin relative to the social allocation.
\end{conditionalresult}

The inequality $Q_t>q_t$ is not derived in the current model. Establishing it
requires a planner costate equation evaluated consistently with the private
allocation and an explicit source of consumer-surplus or knowledge spillovers.

Static market power restricts current AI services. A second distortion arises if an
explicit planner problem implies $Q_t>q_t$, in which case incomplete appropriation
of consumer surplus or knowledge benefits depresses frontier research. Price
regulation may then improve current deployment while weakening the rents that
finance research. The current model establishes the first distortion but leaves the
sign of the research wedge conditional.
